\documentclass[UTF8,a4paper,12pt]{ctexart}

\usepackage[a4paper,margin=2.2cm]{geometry}
\usepackage{array}
\usepackage{booktabs}
\usepackage{enumitem}
\usepackage{fancyhdr}
\usepackage{hyperref}
\usepackage{listings}
\usepackage[table]{xcolor}
\usepackage{tabularx}
\usepackage{tikz}
\usepackage[most]{tcolorbox}
\usepackage{titlesec}
\usepackage{longtable}
\usepackage{multicol}
\usepackage{etoolbox}
\usetikzlibrary{arrows.meta,positioning,shapes.geometric,fit,calc}

\hypersetup{
  colorlinks=true,
  linkcolor=Summary,
  urlcolor=blue!60!black
}

\definecolor{Accent}{HTML}{2F5D50}
\definecolor{AccentLight}{HTML}{E7F1EE}
\definecolor{Warn}{HTML}{A74B32}
\definecolor{WarnLight}{HTML}{FCEEE9}
\definecolor{Summary}{HTML}{385B8A}
\definecolor{SummaryLight}{HTML}{EAF0F8}
\definecolor{CodeBg}{HTML}{F5F7F8}
\definecolor{GrayLine}{HTML}{D7DDDF}
\definecolor{CoverSand}{HTML}{F4E8D8}

\pagestyle{fancy}
\fancyhf{}
\fancyhead[L]{Codex CLI 使用教程}
\fancyfoot[C]{\thepage}
\setlength{\headheight}{15pt}

\titleformat{\section}[display]
  {\Huge\bfseries\color{Accent}}
  {\color{Summary}\rule{\textwidth}{1.2pt}\\[-0.25em]\Large\thesection}
  {0.35em}
  {\filright}
\titlespacing*{\section}{0pt}{0pt}{1.1em}
\titleformat{\subsection}{\large\bfseries\color{Accent}}{\thesubsection}{0.6em}{}
\setlist[itemize]{leftmargin=1.8em}
\setlist[enumerate]{leftmargin=1.8em}
\setlength{\parindent}{2em}
\setlength{\parskip}{0.35em}

\lstdefinestyle{cli}{
  backgroundcolor=\color{CodeBg},
  basicstyle=\ttfamily\small,
  frame=single,
  rulecolor=\color{GrayLine},
  columns=fullflexible,
  breaklines=true,
  breakautoindent=false,
  breakindent=0pt,
  showstringspaces=false
}

\lstdefinestyle{prompt}{
  style=cli,
  linewidth=0.9\linewidth,
  xleftmargin=0.05\linewidth
}

\newtcolorbox{learnbox}{
  colback=AccentLight,
  colframe=Accent,
  boxrule=0.7pt,
  arc=2mm,
  title=这一节学完你能做什么,
  fonttitle=\bfseries
}

\newtcolorbox{tipbox}{
  colback=SummaryLight,
  colframe=Summary,
  boxrule=0.7pt,
  arc=2mm,
  title=提示,
  fonttitle=\bfseries
}

\newtcolorbox{warnbox}{
  colback=WarnLight,
  colframe=Warn,
  boxrule=0.7pt,
  arc=2mm,
  title=注意,
  fonttitle=\bfseries
}

\newtcolorbox{sumbox}{
  colback=AccentLight,
  colframe=Accent,
  boxrule=0.7pt,
  arc=2mm,
  title=小结,
  fonttitle=\bfseries
}

\newtcolorbox{advancedbox}{
  colback=white,
  colframe=black!45,
  boxrule=0.6pt,
  arc=2mm,
  title=进阶阅读,
  fonttitle=\bfseries
}

\newtcolorbox{versionbox}{
  colback=white,
  colframe=Summary,
  boxrule=0.7pt,
  arc=2mm,
  title=版本提示,
  fonttitle=\bfseries
}

\pretocmd{\section}{\clearpage}{}{}

\begin{document}

\begin{titlepage}
\thispagestyle{empty}
\begin{tikzpicture}[remember picture,overlay]
  \fill[Accent!10] ([xshift=-1.8cm,yshift=1.5cm]current page.north east) circle (4.3cm);
  \fill[Summary!12] ([xshift=1.0cm,yshift=-0.8cm]current page.north west) circle (2.6cm);
  \fill[CoverSand] ([xshift=-1.5cm,yshift=-1.2cm]current page.south east) circle (3.8cm);
  \fill[Accent] ([xshift=2.2cm,yshift=-2.1cm]current page.north west) rectangle ++(5.6cm,-0.24cm);
  \fill[Summary] ([xshift=-7.2cm,yshift=-2.55cm]current page.north east) rectangle ++(3.4cm,-0.18cm);
\end{tikzpicture}

\vspace*{2.6cm}

{\Large\bfseries\color{Accent}OpenAI Codex}\\[0.6em]
{\zihao{0}\bfseries Codex CLI 使用教程}\\[0.8em]
{\Large\color{black!72}从第一次上手，到形成稳定工作流}\\[1.2em]

\begin{tcolorbox}[
  enhanced,
  colback=white,
  colframe=Accent,
  boxrule=1pt,
  arc=3mm,
  left=8mm,
  right=8mm,
  top=7mm,
  bottom=7mm
]
\large
这是一份面向初学者的中文教程。它不追求“参数大全”，而是把安装、提问、权限、审改动、自动化和排错串成一条真正能落地的使用路径。
\end{tcolorbox}

\vspace{0.8cm}

\begin{tabularx}{\textwidth}{>{\raggedright\arraybackslash}X >{\raggedright\arraybackslash}X >{\raggedright\arraybackslash}X}
\tcbox[colback=AccentLight,colframe=Accent,arc=2mm,boxrule=0.8pt,left=2mm,right=2mm,top=1mm,bottom=1mm]{\textbf{适合读者}} &
\tcbox[colback=SummaryLight,colframe=Summary,arc=2mm,boxrule=0.8pt,left=2mm,right=2mm,top=1mm,bottom=1mm]{\textbf{阅读重点}} &
\tcbox[colback=WarnLight,colframe=Warn,arc=2mm,boxrule=0.8pt,left=2mm,right=2mm,top=1mm,bottom=1mm]{\textbf{交付目标}} \\
会一点命令行，但第一次系统使用 Codex CLI 的开发者。 &
真实工作流、提示词写法、权限边界、Git 与 CI（持续集成）配合。 &
看完就能独立完成一轮“理解项目、定位问题、做小改动、复核结果”。 \\
\end{tabularx}

\vfill

\begin{tcolorbox}[
  colback=CodeBg,
  colframe=GrayLine,
  boxrule=0.6pt,
  arc=2mm,
  left=6mm,
  right=6mm,
  top=4mm,
  bottom=4mm
]
\small
本文根据本地 \texttt{codex --help} 输出与 OpenAI 官方 Codex 文档整理；请以你当前安装版本为准。
\end{tcolorbox}
\end{titlepage}

\tableofcontents
\newpage

\section*{教程大纲}
\addcontentsline{toc}{section}{教程大纲}
\begin{itemize}
  \item \textbf{第一部分：先把上手主线跑通}\\
  从“Codex CLI 是什么”一路到“最小可运行示例”，先解决安装、登录、进目录、提任务、看权限、看结果这条最短路径。
  \item \textbf{第二部分：把一轮任务做完整}\\
  用三轮对话、提示词模板、典型场景、Git / 测试 / \texttt{exec} 配合、验收卡和排错思路，把“会用”变成“能落地”。
  \item \textbf{第三部分：长期稳定使用}\\
  这里会讲最佳实践、新手易坑、高风险任务边界，以及“从会用到会控场”的进阶技巧。
  \item \textbf{第四部分：速查与附录}\\
  正文后面先给速读总结、命令速查表和避坑清单；再按需查练习、演练页、\texttt{AGENTS.md}、CI、完整命令参考、FAQ 和官方资源。
\end{itemize}

\begin{tipbox}
第一次通读，建议先把前三部分顺着看完；第四部分更像工具箱，遇到对应场景再翻会更高效。
\end{tipbox}

\newpage
\section{Codex CLI 是什么}

\begin{learnbox}
学完这一节，你会知道 Codex CLI 适合干什么、不适合干什么，以及为什么很多人把它当成“终端里的代码搭档”。
\end{learnbox}

先用一句最朴素的话解释：\textbf{Codex CLI 是一个跑在终端里的 AI 编码代理。}你在项目目录里和它对话，它会结合当前代码、目录、命令行环境，帮你读代码、改代码、跑命令、解释报错、做代码审查。

它最适合的事，通常是下面四类：
\begin{itemize}
  \item 接手陌生项目时，先帮你找入口、找主流程、给阅读顺序。
  \item 你知道要改哪里，但不想自己先翻几十个文件。
  \item 你遇到构建失败、测试失败、脚本报错，想先得到一个靠谱的排查顺序。
  \item 你要补文档、补测试、改小功能，希望先出一个可用初稿。
\end{itemize}

它不太适合的事，也要提前说清楚：
\begin{itemize}
  \item 你完全没有目标，只想“让它自动把整个项目做完”。
  \item 你准备把高风险命令、删除操作、线上变更全交给它，不做复核。
  \item 你把需求判断、架构取舍、业务责任也一起外包给它。
\end{itemize}

\subsection{先看一张工作流程图}

\begin{center}
\resizebox{\textwidth}{!}{%
\begin{tikzpicture}[
  node distance=10mm and 7mm,
  box/.style={draw=Accent, rounded corners, fill=AccentLight, minimum width=2.8cm, minimum height=1cm, align=center},
  line/.style={-Latex, thick, draw=Accent}
]
\node[box] (need) {提出任务\\“帮我定位登录逻辑”};
\node[box, right=of need] (context) {读取上下文\\目录、文件、命令、报错};
\node[box, right=of context] (plan) {生成方案\\解释、改动、验证建议};
\node[box, right=of plan] (act) {执行动作\\读文件、改文件、跑命令};
\node[box, below=of plan] (check) {返回结果\\总结、风险、下一步};

\draw[line] (need) -- (context);
\draw[line] (context) -- (plan);
\draw[line] (plan) -- (act);
\draw[line] (act) |- (check);
\draw[line] (check) -| (context);
\end{tikzpicture}
}
\end{center}

\subsection{一个很实际的场景}

比如你打开一个陌生仓库，里面有十几个目录。你现在既不知道入口文件在哪，也不知道应该从前端、后端还是脚本开始看。你可以让 Codex CLI：
\begin{itemize}
  \item 先扫目录结构；
  \item 再指出入口文件和关键目录；
  \item 然后按新手视角给阅读顺序；
  \item 最后在你确认范围后做小改动。
\end{itemize}

\begin{tipbox}
把 Codex CLI 理解成“既能看代码也能动手，但需要你给目标和边界的工程搭档”，这个心智模型最稳。
\end{tipbox}

\begin{sumbox}
Codex CLI 的价值，不是神奇自动化，而是把“理解项目、修改项目、解释结果、验证改动”放到一个终端工作流里。
\end{sumbox}

\section{安装前准备}

\begin{learnbox}
学完这一节，你会知道在安装前至少要确认哪些东西，避免工具装上了，却因为目录、账号或系统环境卡住。
\end{learnbox}

很多新手的问题不是“不会装”，而是“装完以后不知道为什么不能用”。所以安装前先确认下面几件事。

\subsection{你至少要准备什么}

\begin{itemize}
  \item 一台能正常使用终端的电脑。
  \item 一个你愿意拿来练手的项目目录，最好不是最核心的正式仓库。
  \item 可用的认证方式：ChatGPT 账号登录，或者 OpenAI API 密钥。
  \item 基本的命令行习惯，比如会用 \texttt{cd}、\texttt{ls}、\texttt{pwd}。
\end{itemize}

\subsection{安装前建议先确认目录}

\begin{lstlisting}[style=cli]
pwd
ls
\end{lstlisting}

这两条命令看起来很基础，但足够帮你避免很多低级错误：
\begin{itemize}
  \item 先确认你现在是不是站在目标项目里；
  \item 再确认这个目录到底有哪些文件和子目录。
\end{itemize}

\begin{warnbox}
如果你还没搞清楚“当前目录是不是目标项目”，先别急着让 Codex 改东西。这是新手最常见的坑之一。
\end{warnbox}

\subsection{再确认一下你的使用目标}

安装前最好先回答自己一个问题：你现在主要是想用它做什么？
\begin{itemize}
  \item 只是想交互式地看项目、改小功能。
  \item 想把某些固定任务做成脚本或 CI（Continuous Integration，持续集成）。
  \item 想在本地终端里配合 Git 做评审和修复。
\end{itemize}

不同目标会影响你后面更常用的功能。比如日常开发更常用 \texttt{codex}、\texttt{/diff}、\texttt{/review}；脚本化更常用 \texttt{codex exec}。

\begin{sumbox}
安装前先确认三件事：目录、认证方式、使用目标。先把“在哪里用、怎么登、打算拿它做什么”说清楚，后面会顺很多。
\end{sumbox}

\section{安装、登录与首次验证}

\begin{learnbox}
学完这一节，你会知道怎么把 Codex CLI 真的装起来、登录进去，并完成第一次最小验证。
\end{learnbox}

\subsection{先装上}

OpenAI 官方 Codex CLI 页面当前给出了 \texttt{npm} 安装入口。下面用最常见的方式举例：

\begin{lstlisting}[style=cli]
npm i -g @openai/codex
\end{lstlisting}

装完后可以先确认命令是否存在：

\begin{lstlisting}[style=cli]
codex --version
codex --help
\end{lstlisting}

如果 \texttt{--help} 能正常输出，说明至少两件事成立了：
\begin{itemize}
  \item CLI 已经安装成功；
  \item 当前终端能找到 \texttt{codex} 这个命令。
\end{itemize}

\subsection{登录有两条主路}

官方文档当前写得很明确：Codex 支持两种 OpenAI 登录方式。
\begin{itemize}
  \item 用 ChatGPT 登录：适合日常交互式使用，是 CLI 默认的登录路径。
  \item 用 API 密钥登录：适合按 API 计费的使用方式，也常用于自动化和 CI。
\end{itemize}

\subsection{方式一：用 ChatGPT 登录}

最直接的方式通常就是先跑：

\begin{lstlisting}[style=cli]
codex
\end{lstlisting}

第一次运行时，CLI 会提示你登录。按官方文档，默认会走 ChatGPT 登录流程，浏览器会打开登录页，完成后再回到终端。

如果你想显式发起登录，也可以：

\begin{lstlisting}[style=cli]
codex login
codex login status
\end{lstlisting}

\subsection{方式二：用 API 密钥登录}

如果你更希望按 API 使用方式来计费，或者你就是打算后面接脚本，那 API 密钥登录很实用。本地 CLI 帮助里给出了这类写法：

\begin{lstlisting}[style=cli]
printenv OPENAI_API_KEY | codex login --with-api-key
\end{lstlisting}

它的意思很简单：把环境变量里的密钥通过标准输入传给登录命令。

这里顺手区分一下两个常见变量名：\texttt{OPENAI\_API\_KEY} 常出现在“把密钥传给 \texttt{codex login --with-api-key}”这类登录示例里；到了 \texttt{codex exec} 和 CI 场景，官方更常用的是 \texttt{CODEX\_API\_KEY}。

\begin{warnbox}
别把 API 密钥直接写进教程截图、聊天记录或仓库文件里。环境变量是最基本也最稳妥的做法。
\end{warnbox}

\subsection{无浏览器环境怎么办}

如果你在远端机器、容器或者没有图形界面的环境里工作，官方文档给了两条思路：
\begin{itemize}
  \item 设备码登录：运行 \texttt{codex login --device-auth}，按提示在浏览器里完成一次性验证。
  \item 在本地登录后，把认证缓存复制过去。
\end{itemize}

官方文档还特别提醒：如果你使用文件形式缓存认证信息，\texttt{\textasciitilde/.codex/auth.json} 需要当成密码对待，绝对不要提交到仓库。

\subsection{Linux 环境下最常见的三种起步方式}

下面默认你用的是 Linux 终端。对新手来说，最常见的其实就是三种环境：
\begin{itemize}
  \item \textbf{本地图形界面 Linux}：最省心，直接在项目目录里运行 \texttt{codex}，按提示登录即可。
  \item \textbf{远端 Linux / SSH}：通常没有浏览器，优先用 \texttt{codex login --device-auth}，或者先在本地登录，再把认证缓存安全地带过去。
  \item \textbf{容器或云主机里的 Linux}：除了登录，还要确认当前挂载目录是不是你真想操作的工作区，以及容器销毁后认证信息会不会丢。
\end{itemize}

不管是哪一种，开始前都建议先跑一组最小检查：

\begin{lstlisting}[style=cli]
pwd
ls
codex --help
codex login status
\end{lstlisting}

这四步的作用分别很实际：
\begin{itemize}
  \item \texttt{pwd}：确认你现在到底站在哪个目录；
  \item \texttt{ls}：确认眼前是不是你预期的项目；
  \item \texttt{codex --help}：确认 CLI 本身可用；
  \item \texttt{codex login status}：确认认证状态是不是已经就绪。
\end{itemize}

\begin{tipbox}
如果你在远端 Linux 上工作，最容易忽略的不是命令本身，而是“我现在到底在哪个目录、这个目录是不是可写、认证是不是还在”。先把这三件事确认掉，后面会顺很多。
\end{tipbox}

\subsection{macOS、Windows 和 WSL 用户额外看这几眼}

如果你不是在纯 Linux 环境里工作，也完全没问题。真正容易卡住的，不是 Codex 本身，而是终端、路径、认证和项目目录混在一起了。

\begin{itemize}
  \item \textbf{macOS}：整体体验和 Linux 很接近。重点确认当前终端是不是你常用的 shell，\texttt{codex} 能不能被 PATH 找到，以及你打开的是不是正确的项目目录。
  \item \textbf{Windows PowerShell}：命令能跑，但环境变量写法和 Bash 不一样。临时设置 API 密钥时，更像下面这样：
\end{itemize}

\begin{lstlisting}[style=cli]
$env:OPENAI_API_KEY="your_api_key"
$env:CODEX_API_KEY="your_api_key"
\end{lstlisting}

\begin{itemize}
  \item \textbf{WSL}：最好把它当成“一台单独的 Linux 环境”来理解。认证、Git、Node、测试命令，尽量都在同一个 WSL 终端里完成。
\end{itemize}

如果你主要在 WSL 里开发，一个很稳的顺序通常是：
\begin{enumerate}
  \item 把项目放在你真正准备开发的那一侧；
  \item 在那一侧完成登录、运行 \texttt{codex}、跑 Git 和测试；
  \item 不要一会儿在 Windows 终端里改，一会儿又切去 WSL 里验证。
\end{enumerate}

\begin{warnbox}
最容易出问题的场景，是“仓库在 WSL 里，终端却在 Windows 那边；或者认证在一边，Git 和依赖在另一边”。如果项目主要跑在 WSL，就尽量把 Codex、Git、测试也都放在 WSL 里统一完成。
\end{warnbox}

\subsection{第一次验证，不要跳过}

建议按下面顺序做一次最小验证：

\begin{enumerate}
  \item \texttt{codex --help} 能输出；
  \item \texttt{codex login status} 能确认已登录；
  \item 进入一个项目目录后，运行 \texttt{codex}；
  \item 提一个最小任务，看看它能不能正确读取项目结构。
\end{enumerate}

\subsection{升级也很简单}

官方 CLI 页面当前给出的 npm 升级方式是：

\begin{lstlisting}[style=cli]
npm i -g @openai/codex@latest
\end{lstlisting}

\begin{tipbox}
新版本发布会比较频繁。你如果发现教程里的某个选项和本地帮助不一致，优先相信你机器上的 \texttt{codex --help}。
\end{tipbox}

\begin{sumbox}
这一章最重要的不是“背命令”，而是形成一个稳定顺序：安装、登录、验证。只要这三步走通，后面的使用成本会大幅下降。
\end{sumbox}

\section{基本使用流程}

\begin{learnbox}
学完这一节，你会知道第一次真正使用 Codex CLI 的标准节奏，至少能独立完成一轮“提任务 -> 看结果 -> 继续追问 -> 做验证”。
\end{learnbox}

对新手来说，最容易上手的不是复杂参数，而是一个简单的工作流：

\begin{enumerate}
  \item 进入项目目录；
  \item 用一句清楚的话描述任务；
  \item 先让它读上下文；
  \item 看它的解释、计划或改动建议；
  \item 不满意就继续追问；
  \item 改完以后自己做验证。
\end{enumerate}

\subsection{典型使用过程步骤图}

\begin{center}
\begin{tikzpicture}[
  stepbox/.style={draw=Summary, rounded corners, fill=SummaryLight, minimum width=3.2cm, minimum height=0.95cm, align=center},
  arr/.style={-Latex, thick, draw=Summary}
]
\node[stepbox] (s1) {1. 进入项目目录};
\node[stepbox, right=6mm of s1] (s2) {2. 说清楚任务};
\node[stepbox, right=6mm of s2] (s3) {3. 让它先看代码};
\node[stepbox, below=9mm of s2] (s4) {4. 查看解释或修改};
\node[stepbox, left=6mm of s4] (s5) {5. 继续追问};
\node[stepbox, right=6mm of s4] (s6) {6. 验证结果};

\draw[arr] (s1) -- (s2);
\draw[arr] (s2) -- (s3);
\draw[arr] (s3) -- (s4);
\draw[arr] (s4) -- (s5);
\draw[arr] (s4) -- (s6);
\end{tikzpicture}
\end{center}

\subsection{第一次提问，建议这样说}

比起一句“帮我看看这个项目”，下面这种说法通常更有效：

\begin{lstlisting}[style=prompt]
请先看一下当前项目的目录结构，告诉我入口文件、主要模块，以及新手应该先看哪几个文件。
\end{lstlisting}

\subsection{追问时，尽量把范围收紧}

如果第一轮结果太泛，你可以继续这样问：

\begin{lstlisting}[style=prompt]
请不要泛泛介绍，重点解释登录相关代码在哪、数据流怎么走，以及我应该先从哪个文件开始读。
\end{lstlisting}

\begin{warnbox}
“帮我优化一下”“帮我改好”这种说法太宽，容易得到看起来很忙、但不够落地的结果。新手阶段尽量把目标说具体。
\end{warnbox}

\begin{sumbox}
第一次用 Codex CLI，不求一步到位。先学会把任务说清楚，再学会继续追问，这比背一堆参数更重要。
\end{sumbox}

\section{常见命令与操作方式}

\begin{learnbox}
学完这一节，你会知道日常最常用的命令有哪些，各自适合什么场景。
\end{learnbox}

\subsection{先看一张分类图}

\begin{center}
\begin{tikzpicture}[
  main/.style={draw=Accent, rounded corners, fill=AccentLight, minimum width=3.4cm, minimum height=1cm, align=center},
  sub/.style={draw=black!50, rounded corners, fill=white, minimum width=3.15cm, minimum height=0.85cm, align=center},
  line/.style={-Latex, thick}
]
\node[main] (m) {Codex CLI 常见操作};
\node[sub, below left=10mm and 16mm of m] (a) {交互式使用\\\texttt{codex}};
\node[sub, below=16mm of m] (b) {脚本化执行\\\texttt{codex exec}};
\node[sub, below right=10mm and 16mm of m] (c) {会话管理\\\texttt{resume / fork}};
\node[sub, below=10mm of a] (d) {登录与帮助};
\node[sub, below=10mm of b] (e) {代码评审与输出控制};
\node[sub, below=10mm of c] (f) {恢复会话与应用改动};

\draw[line] (m) -- (a);
\draw[line] (m) -- (b);
\draw[line] (m) -- (c);
\draw[line] (a) -- (d);
\draw[line] (b) -- (e);
\draw[line] (c) -- (f);
\end{tikzpicture}
\end{center}

\subsection{交互式入口：\texttt{codex}}

\begin{lstlisting}[style=cli]
codex
\end{lstlisting}

这是最常见的入口。适合下面这些场景：
\begin{itemize}
  \item 你要围绕当前项目连续提问；
  \item 你想边看解释边调整方向；
  \item 你希望它先分析，再改动，再复盘。
\end{itemize}

\subsection{最基础的帮助与状态}

\begin{lstlisting}[style=cli]
codex --help
codex --version
codex login status
\end{lstlisting}

忘了支持什么功能、当前版本是什么、是否已经登录，先看这些，比乱猜靠谱。

\subsection{非交互执行：\texttt{codex exec}}

\begin{lstlisting}[style=cli]
codex exec "检查当前仓库里最可能导致测试失败的改动"
\end{lstlisting}

这类命令更适合脚本、流水线、一次性任务。官方文档也明确把 \texttt{codex exec} 定位为“脚本化或 CI 风格的非交互执行模式”。

\subsection{代码评审：\texttt{codex review}}

\begin{lstlisting}[style=cli]
codex review --uncommitted
codex review --base main
\end{lstlisting}

从本地帮助看，\texttt{review} 适合两类场景：
\begin{itemize}
  \item 你还没提交，只想先看当前未提交改动的问题；
  \item 你想把当前分支和某个基线分支做对比审查。
\end{itemize}

\subsection{恢复、分叉、应用结果}

\begin{lstlisting}[style=cli]
codex resume --last
codex fork --last
codex apply <TASK_ID>
\end{lstlisting}

这几个命令很实用：
\begin{itemize}
  \item \texttt{resume}：接着上次会话干，不用重新讲背景。
  \item \texttt{fork}：把当前思路分成两条线，试不同方案。
  \item \texttt{apply}：把某次任务产生的 diff 应到本地工作区。
\end{itemize}

\subsection{两个高频参数}

\begin{lstlisting}[style=cli]
codex --search
codex --add-dir ../shared-lib
\end{lstlisting}

\texttt{--search} 会打开实时网页搜索能力；\texttt{--add-dir} 用来在主工作目录之外，再给 Codex 多一个可写目录。

\subsection{什么时候该开 \texttt{--search}，什么时候只看本地代码}

\texttt{--search} 很有用，但不适合默认一直开着。最稳的判断方法很简单：\textbf{只要问题的答案主要藏在当前仓库里，就先别联网；只有当你真的需要外部最新资料时，再把搜索打开。}

更适合\textbf{只看本地代码}的情况，通常有这些：
\begin{itemize}
  \item 你要找入口文件、调用链、核心目录；
  \item 你要判断当前分支到底改了什么；
  \item 你要分析仓库里已有的报错、测试、脚本和配置；
  \item 你要决定“最小改动点”到底落在哪个文件。
\end{itemize}

更适合\textbf{打开 \texttt{--search}} 的情况，通常有这些：
\begin{itemize}
  \item 你要查官方文档、SDK 当前版本行为、参数变化；
  \item 你要看外部 API、依赖库、平台规则的最新说明；
  \item 你要确认某个错误码、版本差异、产品限制是不是最近才变过；
  \item 你明确需要来源链接，而不是只要一句经验判断。
\end{itemize}

一个很稳的问法，是把本地事实和外部资料拆成两轮：

\begin{lstlisting}[style=prompt]
先只基于当前仓库回答：登录失败最相关的文件、调用链和最小改动点是什么？
\end{lstlisting}

\begin{lstlisting}[style=prompt]
再结合官方文档或可信来源，补充说明这里依赖的 SDK 在当前版本里有没有已知限制。
请把“仓库里实际行为”和“外部资料里的说法”分开写，并给出来源。
\end{lstlisting}

\begin{warnbox}
如果你一上来就把联网搜索和本地代码问题混在一起问，很容易得到“看起来很懂，但落不到当前仓库”的回答。先拿到本地事实，再补外部资料，通常更稳。
\end{warnbox}

\begin{tipbox}
如果你只记三个入口，先记 \texttt{codex}、\texttt{codex exec}、\texttt{codex review}。这三类已经覆盖了大多数入门场景。
\end{tipbox}

\begin{sumbox}
这一章的重点不是全部记住，而是知道怎么分：交互式用 \texttt{codex}，脚本化用 \texttt{codex exec}，改动复核用 \texttt{codex review}。
\end{sumbox}

\section{权限、沙箱与审批机制}

\begin{learnbox}
学完这一节，你会明白为什么 Codex 有时能直接跑命令，有时会停下来问你；也会知道不同沙箱模式分别意味着什么。
\end{learnbox}

这部分很重要。很多人第一次用 Codex CLI，会困惑于下面这些现象：
\begin{itemize}
  \item 为什么它能读文件，但不能写？
  \item 为什么有些命令会要求你批准？
  \item 为什么默认情况下网络能力受限？
  \item 为什么 CI 里和本地交互时的习惯不一样？
\end{itemize}

答案基本都和\textbf{沙箱（sandbox）}和\textbf{审批策略（approval policy）}有关。

\subsection{先记住三个沙箱级别}

\rowcolors{2}{white}{AccentLight}
\begin{longtable}{p{0.24\textwidth} p{0.18\textwidth} p{0.42\textwidth}}
\toprule
模式 & 大致含义 & 适合什么场景 \\
\midrule
\endhead
\texttt{read-only} & 只能读，不能改 & 你只是想看项目、看报错、做分析 \\
\texttt{workspace-write} & 可以在工作区内读写 & 本地日常开发，最常用 \\
\texttt{danger-full-access} & 权限非常大 & 只在你非常明确风险时才考虑 \\
\bottomrule
\end{longtable}

\subsection{再记住三个审批策略}

\begin{itemize}
  \item \texttt{untrusted}：只让它自动跑已知安全的读操作，不信任的命令先问你。
  \item \texttt{on-request}：由代理判断什么时候需要向你请求批准。
  \item \texttt{never}：不弹审批，适合自动化场景，但要配合你自己选好的沙箱。
\end{itemize}

\subsection{最常见的组合}

官方文档里给了几组很实用的默认组合，可以直接拿来理解：

\begin{itemize}
  \item 自动但相对稳妥：
\end{itemize}

\begin{lstlisting}[style=cli]
--full-auto
\end{lstlisting}

它本质上等于：

\begin{lstlisting}[style=cli]
--sandbox workspace-write --ask-for-approval on-request
\end{lstlisting}

\begin{itemize}
  \item 只读浏览：
\end{itemize}

\begin{lstlisting}[style=cli]
--sandbox read-only --ask-for-approval on-request
\end{lstlisting}

\begin{itemize}
  \item 工作区可写，但不信任命令必须先批：
\end{itemize}

\begin{lstlisting}[style=cli]
--sandbox workspace-write --ask-for-approval untrusted
\end{lstlisting}

\begin{itemize}
  \item 完全放开：
\end{itemize}

\begin{lstlisting}[style=cli]
--dangerously-bypass-approvals-and-sandbox
\end{lstlisting}

\subsection{为什么很多人推荐从 \texttt{workspace-write} 开始}

因为它比较符合本地开发的真实需求：
\begin{itemize}
  \item 能改你当前项目里的文件；
  \item 不会默认把整个机器都放开；
  \item 和审批策略配合时，心里更有数。
\end{itemize}

\subsection{什么时候会要求你批准}

官方文档明确提到，在 \texttt{untrusted} 模式下，Codex 会自动运行已知安全的读操作；但\textbf{会改变状态、可能有副作用、或会走外部执行路径的命令}，通常都要你批准。比如某些破坏性 Git 操作，就是典型高风险操作。

\begin{warnbox}
\texttt{--dangerously-bypass-approvals-and-sandbox} 这个名字已经说得很直白了。只有在你外层环境本身就已经被严格隔离时，才考虑用它。
\end{warnbox}

\subsection{一个判断思路图}

\begin{center}
\begin{tikzpicture}[
  q/.style={diamond, draw=Warn, fill=WarnLight, aspect=2, align=center, inner sep=2pt},
  r/.style={draw=black!55, rounded corners, fill=white, minimum width=3.6cm, minimum height=0.9cm, align=center},
  arr/.style={-Latex, thick, draw=Warn}
]
\node[q] (start) {你现在主要想做什么};
\node[r, below left=10mm and 14mm of start] (read) {只读项目\\先用 read-only};
\node[r, below=16mm of start] (dev) {本地日常开发\\先用 workspace-write};
\node[r, below right=10mm and 14mm of start] (auto) {自动化 / CI\\考虑 never + 明确沙箱};
\node[r, below=14mm of dev] (risk) {高风险操作前\\自己复核并批准};

\draw[arr] (start) -- (read);
\draw[arr] (start) -- (dev);
\draw[arr] (start) -- (auto);
\draw[arr] (read) |- (risk);
\draw[arr] (dev) -- (risk);
\draw[arr] (auto) |- (risk);
\end{tikzpicture}
\end{center}

\begin{sumbox}
最实用的经验是：本地开发从 \texttt{workspace-write} 起步，先别碰完全放开模式；你越清楚权限边界，Codex 用起来就越稳。
\end{sumbox}

\section{实用斜杠命令与会话控制}

\begin{learnbox}
学完这一节，你会知道在交互会话里有哪些高频斜杠命令，不用退出重进也能切模型、看状态、看 diff、做评审。
\end{learnbox}

CLI 里很多高频操作，不一定要回到命令行里敲完整命令。官方文档的斜杠命令页里，有几条特别适合新手先记住。

\subsection{先看这张入口选择图}

\begin{center}
\resizebox{\textwidth}{!}{%
\begin{tikzpicture}[
  q/.style={draw=Summary, diamond, aspect=2.1, fill=SummaryLight, align=center, inner sep=2.5pt},
  a/.style={draw=Accent, rounded corners, fill=AccentLight, minimum width=3.4cm, minimum height=0.95cm, align=center},
  line/.style={-Latex, thick}
]
\node[q] (start) {你现在想做什么？};
\node[a, below left=14mm and 34mm of start] (chat) {连续对话、逐步迭代\\用 \texttt{codex}};
\node[a, below right=14mm and 34mm of start] (exec) {脚本、CI、一次性任务\\用 \texttt{codex exec}};
\node[q, below=16mm of chat] (inside) {已经在会话里了，还想继续怎么走？};
\node[a, below left=14mm and 28mm of inside] (new) {当前终端里重开一轮\\\texttt{/new} 或 \texttt{/clear}};
\node[a, below=14mm of inside] (resume) {切回之前保存的会话\\\texttt{/resume} 或 \texttt{codex resume --last}};
\node[a, below right=14mm and 28mm of inside] (fork) {保留旧讨论再开新线\\\texttt{/fork} 或 \texttt{codex fork --last}};

\draw[line] (start) -- (chat);
\draw[line] (start) -- (exec);
\draw[line] (chat) -- (inside);
\draw[line] (inside) -- (new);
\draw[line] (inside) -- (resume);
\draw[line] (inside) -- (fork);
\end{tikzpicture}
}
\end{center}

\begin{tipbox}
一句话判断：要连续聊就进 \texttt{codex}；要脚本化就用 \texttt{codex exec}；已经在会话里还想换路，就优先想 \texttt{/new}、\texttt{/resume}、\texttt{/fork}。
\end{tipbox}

\subsection{最值得先记的几条}

\rowcolors{2}{white}{SummaryLight}
\begin{longtable}{>{\ttfamily}p{0.18\textwidth} p{0.22\textwidth} p{0.42\textwidth}}
\toprule
斜杠命令 & 作用 & 最常见用途 \\
\midrule
\endhead
/status & 看当前会话状态 & 确认模型、审批策略、可写目录、token 使用情况 \\
/model & 切模型 & 当前任务需要更快或更深的推理时切换 \\
/permissions & 改权限模式 & 临时把会话改成更保守或更自动 \\
/plan & 先出执行计划 & 任务较大时先计划再动手 \\
/diff & 看 Git diff & 改完先审一眼再提交 \\
/review & 做工作区评审 & 查行为风险、缺测试、潜在回归 \\
/new & 开一段新会话 & 想在当前 CLI 里重开一轮，不退出终端 \\
/clear & 清屏并新开会话 & 想把终端内容清掉，再开始新对话 \\
/resume & 恢复旧会话 & 想在交互里切到之前保存的会话 \\
/fork & 从旧会话分叉 & 想保留旧上下文，但另开一条新线 \\
/mention & 显式提某个文件 & 强制把某个路径拉进上下文 \\
/compact & 压缩长会话 & 对话很长时瘦身，保留重点 \\
/init & 生成 \texttt{AGENTS.md} 脚手架 & 把仓库规则写成长期提示 \\
\bottomrule
\end{longtable}

\subsection{\texttt{/status}：先确认它到底在什么状态}

官方文档对 \texttt{/status} 的描述很实用：它会显示当前模型、审批策略、可写根目录、token 使用情况。也就是说，当你觉得“它怎么突然行为不对劲”时，先跑 \texttt{/status} 通常最省时间。

\subsection{\texttt{/permissions}：会话里随时改权限}

你一开始可以保守一些，先读代码；真正要改时，再把权限调到更适合当前任务的模式。这样比整场会话从头到尾都开得很大更稳。

\subsection{\texttt{/review} 和 \texttt{/diff}：改完一定要连着看}

这两条更稳的顺序是连着用：
\begin{enumerate}
  \item 先用 \texttt{/review} 让 Codex 从行为和测试角度看一下工作区；
  \item 再用 \texttt{/diff} 看看具体改了哪些文件、哪些行。
\end{enumerate}

\subsection{\texttt{/plan} 和 \texttt{/init}：一个管当前任务，一个管长期规则}

\texttt{/plan} 适合当前任务比较大时先做执行计划。\texttt{/init} 则会在当前目录生成 \texttt{AGENTS.md} 脚手架，这个文件可以把“如何跑项目、如何测试、哪些目录别碰、完成标准是什么”这些长期规则写进去。

\begin{tipbox}
如果你经常重复说“先跑测试再改”“不要动生成文件”“改完给我验证步骤”，就该考虑用 \texttt{/init} 和 \texttt{AGENTS.md} 固化这些规则了。
\end{tipbox}

\subsection{\texttt{/compact} 和 \texttt{/mention}：上下文太长时特别有用}

\texttt{/compact} 适合会话已经很长、背景很多，但你还想继续沿着同一个目标聊下去的时候。它的核心作用是把长对话压缩成更短的可继续上下文。\\
\texttt{/mention} 则适合你想强制把某个文件、目录或路径拉进当前上下文，不希望 Codex 只做泛泛回答的时候。

\begin{tipbox}
如果你已经开始反复贴长日志、长代码、长历史结论，就该优先考虑 \texttt{/compact}；如果它总是没有抓到你要看的文件，就用 \texttt{/mention} 点名。
\end{tipbox}

\subsection{\texttt{/model}：什么时候值得切模型}

\texttt{/model} 适合两种情况：
\begin{itemize}
  \item 你想先用更快的模型做浏览、定位、整理；
  \item 你准备进入更复杂的分析、修改或审查，想换到更强的模型。
\end{itemize}

\subsection{\texttt{/new} 和 \texttt{/clear}：在当前终端里重开一轮}

\texttt{/new} 就是用来在\textbf{当前 CLI 里开启一段新的对话} 的，不用先退出终端。\\
\texttt{/clear} 的作用更直接一些：\textbf{先清空当前屏幕显示，再开始一段新的对话}。如果你只是想重新聊一轮，但不在乎清不清屏，用 \texttt{/new} 就够了。

\begin{lstlisting}[style=cli]
/new
/clear
\end{lstlisting}

\subsection{\texttt{/resume} 和 \texttt{/fork}：恢复旧会话，或者从旧会话分叉}

\texttt{/resume} 用来恢复之前保存过的会话；\texttt{/fork} 则是基于已有会话分一条新线。对于“我能不能在多个会话间切换”这个问题，最接近的答案是：\textbf{可以切到别的已保存会话，但方式是恢复或分叉，不是像浏览器标签页那样常驻并排切。}

\begin{lstlisting}[style=cli]
/resume
/fork
\end{lstlisting}

如果你在命令行里操作，也可以直接用：

\begin{lstlisting}[style=cli]
codex resume --last
codex fork --last
\end{lstlisting}

\subsection{什么时候退出后重开 \texttt{codex}？}

如果你是想连终端态都一起重置，或者你已经退出了交互，那当然也可以直接重新运行：

\begin{lstlisting}[style=cli]
codex
\end{lstlisting}

但从日常使用角度看，优先级通常是：
\begin{itemize}
  \item 想原地开启一轮新对话：\texttt{/new}
  \item 想清屏后再开：\texttt{/clear}
  \item 想切到旧会话：\texttt{/resume}
  \item 想从旧会话分一条新线：\texttt{/fork}
  \item 想从命令行直接恢复最近一次：\texttt{codex resume --last}
\end{itemize}

\subsection{会话命令对照图}

\begin{center}
\begin{tcbraster}[raster columns=2, raster equal height=rows, raster column skip=6pt, raster row skip=6pt]
\begin{tcolorbox}[colback=SummaryLight, colframe=Summary, title=\texttt{/new}]
在当前终端里直接开一段新的对话。\\
\textbf{适合：}想重开一轮，但不想退出 Codex。
\end{tcolorbox}
\begin{tcolorbox}[colback=AccentLight, colframe=Accent, title=\texttt{/clear}]
先清空当前屏幕，再开始新的对话。\\
\textbf{适合：}内容太乱，想顺手清屏。
\end{tcolorbox}
\begin{tcolorbox}[colback=SummaryLight, colframe=Summary, title=\texttt{/resume}]
恢复之前保存过的会话。\\
\textbf{适合：}切回某个旧会话，继续原来的上下文。
\end{tcolorbox}
\begin{tcolorbox}[colback=AccentLight, colframe=Accent, title=\texttt{/fork}]
基于已有会话分一条新线。\\
\textbf{适合：}保留旧思路，但另开一个方向继续试。
\end{tcolorbox}
\end{tcbraster}
\end{center}

\begin{tipbox}
一句话记忆：\texttt{/new} 是重开，\texttt{/clear} 是清屏重开，\texttt{/resume} 是回到旧会话，\texttt{/fork} 是从旧会话分叉。
\end{tipbox}

\begin{sumbox}
斜杠命令的核心价值，是让你在同一场会话里完成“看状态、改权限、切模型、看 diff、做代码审查”这些高频动作，而不用来回跳出会话。
\end{sumbox}

\section{最小可运行示例}

\begin{learnbox}
学完这一节，你至少能照着完成一轮最小体验：进入目录、发起任务、继续追问、复核结果。
\end{learnbox}

这一节只做一件事：帮你把第一轮真的跑通。它和后面的“三轮对话示例”分工不同，前者是最短路径，后者是教你怎么把同一个问题越问越准。

\subsection{示例目标}

假设你刚进入一个项目目录，想知道：
\begin{itemize}
  \item 这个项目的入口在哪；
  \item 哪几个文件最值得先看；
  \item 如果你是新手，阅读顺序应该是什么。
\end{itemize}

\subsection{一轮最小示例}

\begin{lstlisting}[style=cli]
cd your-project
codex
\end{lstlisting}

进入交互后，你可以这样说：

\begin{lstlisting}[style=prompt]
请先查看当前项目结构，告诉我入口文件、核心模块、建议的阅读顺序，并尽量用新手能听懂的话解释。
\end{lstlisting}

如果结果不错，继续追问：

\begin{lstlisting}[style=prompt]
请进一步解释最核心的两个文件各自负责什么、它们之间怎么配合，以及我如果要改登录逻辑应该先看哪里。
\end{lstlisting}

\subsection{最后别忘了做一次会话内检查}

\begin{lstlisting}[style=cli]
/status
\end{lstlisting}

看看当前模型、权限模式、可写目录是不是和你预期一致。如果要审改动，再接：

\begin{lstlisting}[style=cli]
/diff
\end{lstlisting}

\begin{warnbox}
如果它只给你空泛总结，没有落到具体文件名、模块名和路径，那就继续追问，不要直接收工。
\end{warnbox}

\begin{sumbox}
最小示例的目标不是展示高级能力，而是让你真正建立起“我已经会跑通第一轮”的感觉。
\end{sumbox}

\section{三轮对话示例：怎么把问题越问越准}

\begin{learnbox}
学完这一节，你会知道为什么很多好结果不是“第一问神准”，而是靠两三轮追问把范围一点点压出来。
\end{learnbox}

前一节解决的是“能不能顺利跑起来”，这一节解决的是“同一个问题怎么越问越准”。如果你现在只想先把工具跑通，这一节可以先略读，等真正开始多轮迭代时再回来细看。

很多新手第一轮就想拿到最终答案，但真实工作里，更稳的顺序通常是：先拿地图，再抓主线，最后要求可交付结果。下面用“看懂陌生仓库里的登录逻辑”为例。

\subsection{第一轮：先把地图拿到手}

\begin{lstlisting}[style=prompt]
请先按新手视角查看当前仓库，告诉我入口文件、和登录最相关的目录，以及推荐阅读顺序。
先不要修改代码。
\end{lstlisting}

这一轮的目标不是立刻修东西，而是先让 Codex 把范围从“整个仓库”收成“和登录相关的那几个入口、目录、文件”。

\subsection{第二轮：把范围压到功能级}

\begin{lstlisting}[style=prompt]
现在只聚焦登录流程。请列出最相关的 3 到 5 个文件，按调用顺序解释它们分别做什么，
并指出如果我要改“登录失败提示文案”，最可能先看哪一个文件。
\end{lstlisting}

这一轮你已经不是在听泛泛介绍，而是在逼它回答三件更具体的事：
\begin{itemize}
  \item 到底是哪些文件最相关；
  \item 这些文件之间怎么串起来；
  \item 如果要改一个很小的点，入口最可能在哪。
\end{itemize}

\subsection{第三轮：要求它给可交付结果}

\begin{lstlisting}[style=prompt]
请只修改登录失败提示文案，不要改业务逻辑，也不要顺手重构。
改完后请按“改动文件 / 为什么改 / 如何验证 / 剩余风险”四项总结。
\end{lstlisting}

到这一轮，任务边界已经很清楚了：功能只看登录失败提示，改动只碰文案，不做顺手优化，最后还要交付可审的总结格式。

\begin{warnbox}
如果第二轮你还没拿到具体文件和调用链，就先别急着进入第三轮。追问不是浪费时间，而是在给后面的修改省返工。
\end{warnbox}

\begin{sumbox}
一条很稳的追问节奏是：第一轮拿地图，第二轮抓主线，第三轮要交付。这样通常比“一句话把所有要求塞满”更容易得到靠谱结果。
\end{sumbox}

\section{高质量提示词模板库}

\begin{learnbox}
学完这一节，你会拿到一组可以直接复制、稍改就能用的模板，而不是每次临场现编。
\end{learnbox}

很多时候，Codex 用得好不好，不在于你是不是写了“高级提示词”，而在于你有没有说清楚\textbf{目标、范围、约束、输出格式}。

\subsection{模板一：看懂陌生项目}

\begin{lstlisting}[style=prompt]
请先从新手视角介绍当前项目。
重点告诉我：
1. 入口文件在哪里
2. 核心目录分别负责什么
3. 主流程怎么走
4. 建议我按什么顺序阅读
请尽量引用具体文件路径，不要泛泛而谈。
\end{lstlisting}

\subsection{模板二：定位某个功能}

\begin{lstlisting}[style=prompt]
请帮我定位“登录”相关代码。
我想知道：
1. 入口在哪里
2. 涉及哪些文件
3. 调用链大概怎么走
4. 如果我要改登录失败提示，应该先看哪里
\end{lstlisting}

\subsection{模板三：修一个明确的问题}

\begin{lstlisting}[style=prompt]
我遇到了这个报错，请先不要急着改代码。
请按下面顺序回答：
1. 最可能的原因
2. 建议的排查顺序
3. 最值得先看的文件
4. 如果要修，最小改动方案是什么
\end{lstlisting}

\subsection{模板四：做一个小改动，但限制范围}

\begin{lstlisting}[style=prompt]
请只修改提示文案，不要改业务逻辑。
修改前先告诉我你准备改哪些文件。
修改后请总结：
1. 改了哪些文件
2. 每个文件改了什么
3. 我应该怎么验证
\end{lstlisting}

\subsection{模板五：做代码审查}

\begin{lstlisting}[style=prompt]
请像代码审查者一样检查当前改动。
重点关注：
1. 行为变化
2. 潜在回归
3. 边界情况
4. 缺失的测试
请先给发现的问题，再给简短总结。
\end{lstlisting}

\subsection{模板六：让它先采访你}

这是官方最佳实践页面里很值得学的一条思路：如果你的需求还很模糊，可以先让 Codex 反过来问你问题，把模糊目标问清楚。

\begin{lstlisting}[style=prompt]
我现在只有一个模糊想法，请先通过提问帮我把需求收紧，挑战我的假设，等目标清楚以后再开始修改代码。
\end{lstlisting}

\begin{tipbox}
真正高质量的提示词，不是写得花，而是让 Codex 明白“你到底要它做什么、不许它做什么、最后要怎么交付结果”。
\end{tipbox}

\begin{sumbox}
写提示词时，尽量补齐四要素：目标、范围、限制、输出格式。这个习惯比背模板更重要。
\end{sumbox}

\section{几个典型使用场景}

\begin{learnbox}
学完这一节，你会更容易判断什么时候值得用 Codex CLI，什么时候自己手动做反而更快。
\end{learnbox}

这一节更像场景索引，不是新的主线章节。你不一定要逐段精读；更适合在真正碰到“接手项目”“处理报错”“做代码审查”这些任务时，回来按场景抄一句起手式。

\subsection{场景一：接手陌生项目}

你最需要的不是“全量介绍”，而是“给我一个阅读入口”。

\begin{lstlisting}[style=prompt]
请按新手视角介绍这个项目，重点说入口、核心目录、主要数据流，并给出推荐阅读顺序。
\end{lstlisting}

\subsection{场景二：定位某个功能}

\begin{lstlisting}[style=prompt]
请帮我定位上传功能相关代码，告诉我涉及哪些文件、入口在哪里、调用链大概是什么。
\end{lstlisting}

\subsection{场景三：解释报错}

\begin{lstlisting}[style=prompt]
这是我刚遇到的报错。请先解释最可能的原因，再给我一个排查顺序，最后告诉我应该先看哪些文件。
\end{lstlisting}

\subsection{场景四：改一个小功能}

\begin{lstlisting}[style=prompt]
请把这里的提示文案改得更清楚，但不要改业务逻辑。修改后请说明改了哪些文件，以及如何验证。
\end{lstlisting}

\subsection{场景五：提交前做一次代码审查}

\begin{lstlisting}[style=cli]
/review
/diff
\end{lstlisting}

这是很典型的一套动作：先让它帮你找风险，再自己看差异细节。

\begin{sumbox}
Codex CLI 最值钱的地方，往往不是帮你“一口气写完整个功能”，而是大幅缩短“理解、定位、修改、复核”这几步之间的时间。
\end{sumbox}

\section{和 Git、测试、非交互模式的配合}

\begin{learnbox}
学完这一节，你会知道如何把 Codex CLI 接进真实开发流程，而不是只停留在聊天式体验。
\end{learnbox}

\subsection{和 Git 的最小配合方式}

你可以把一轮工作想成这样：
\begin{enumerate}
  \item 先让 Codex 解释或修改；
  \item 再让它做 \texttt{/review} 或 \texttt{codex review}，先把行为风险和缺失测试圈出来；
  \item 然后用 \texttt{/diff} 或 Git diff 看差异，确认具体改了哪些文件、哪些行；
  \item 最后自己跑测试、决定是否提交。
\end{enumerate}

如果你想用命令行方式做审查，本地帮助里有这些入口：

\begin{lstlisting}[style=cli]
codex review --uncommitted
codex review --base main
\end{lstlisting}

\subsection{改坏了怎么办：先保留现场，再决定怎么撤}

新手最需要补的一课，其实不是“怎么让它改”，而是“改得不对时怎么稳稳收回来”。最怕的不是改坏，而是你在没看清现场时又继续乱改。

更稳的顺序通常是：
\begin{enumerate}
  \item 先看 \texttt{git status}；
  \item 再看 \texttt{git diff}；
  \item 如果当前现场还值得保留，先切一条救援分支，或者先做一次 \texttt{stash}；
  \item 最后才决定是继续修、部分撤回，还是整批撤回。
\end{enumerate}

\begin{lstlisting}[style=cli]
git status
git diff
git switch -c backup/codex-rollback-check
git stash push -u -m "codex-before-rollback"
\end{lstlisting}

如果你只是想撤掉某一个文件的工作区改动，可以先从最小动作开始：

\begin{lstlisting}[style=cli]
git restore path/to/file
\end{lstlisting}

如果你非常确定这批跟踪文件里的改动都不要了，才考虑：

\begin{lstlisting}[style=cli]
git restore .
\end{lstlisting}

这时也可以先让 Codex 帮你“做判断”，而不是“直接执行”：

\begin{lstlisting}[style=prompt]
先不要直接撤回任何东西。请先根据当前 diff 告诉我：
1. 哪些改动明显是这次任务引入的
2. 哪些文件值得保留
3. 如果我要最小回滚，应该先撤哪几个文件
4. 每一步回滚后该怎么验证
\end{lstlisting}

\begin{tipbox}
回滚不等于“一键清空”。先保留证据，再缩小范围，再撤最小那一块，通常比一上来整批回退更稳。
\end{tipbox}

\begin{warnbox}
别把 \texttt{git reset --hard} 这类命令当默认答案。只要你还没想清楚哪些内容要保留，就先别碰这类大锤。
\end{warnbox}

\subsection{改完以后，别只看它说了什么}

真正重要的是验证：
\begin{itemize}
  \item 它改的文件是不是你预期的文件；
  \item 它有没有顺手碰到不该碰的模块；
  \item 构建、测试、关键路径是不是还过得去。
\end{itemize}

\subsection{如果你知道测试命令，最好明确告诉它}

官方最佳实践里有一个很关键的提醒：如果你不告诉 Codex 该怎么构建、怎么测试，它就不容易验证自己的工作。所以你可以在提示里直说：

这里顺手解释一下：\texttt{lint} 一般指\textbf{静态检查}，主要用来提前发现格式、风格和一些明显的代码问题。

\begin{lstlisting}[style=prompt]
改完以后请运行项目已有的测试命令。如果你不确定，先搜索仓库里现有的测试、构建和静态检查脚本再决定。
\end{lstlisting}

\subsection{把长期规则写进 \texttt{AGENTS.md}}

官方最佳实践明确建议，把重复出现的规则沉淀到 \texttt{AGENTS.md}，而不是每次都重新说一遍。一个好用的 \texttt{AGENTS.md} 往往会包括：
\begin{itemize}
  \item 仓库结构和重点目录；
  \item 构建、测试、静态检查命令；
  \item 工程约定和禁止事项；
  \item “算完成”的标准和验证方式。
\end{itemize}

你可以先用：

\begin{lstlisting}[style=cli]
/init
\end{lstlisting}

生成脚手架，再按团队真实习惯补充内容。

\subsection{除了 \texttt{AGENTS.md}，也值得把常用设置固化进 \texttt{config.toml}}

\texttt{AGENTS.md} 管的是\textbf{仓库规则}；\texttt{\textasciitilde/.codex/config.toml} 更适合管你自己的\textbf{默认工作习惯}，比如常用模型、审批策略、沙箱模式，以及几种固定的 profile。

如果你每次进会话都在重复同一种组合，通常就值得把它固化下来。一个很实用的分法是：
\begin{itemize}
  \item \textbf{safe}：只读浏览，适合看项目、看报错、做方案比较；
  \item \textbf{dev}：工作区可写，适合本地日常开发；
  \item \textbf{local}：走本地开源模型，适合做低风险浏览、试验或离线场景。
\end{itemize}

日常使用时，你真正需要记住的，常常只是这些入口：

\begin{lstlisting}[style=cli]
codex -p safe
codex -p dev
codex --oss --local-provider ollama -p local
\end{lstlisting}

如果只是偶尔临时换一次，也不一定要改配置文件，直接在命令行里覆盖就够了：

\begin{lstlisting}[style=cli]
codex -s read-only -a on-request
codex -s workspace-write -a on-request --search
\end{lstlisting}

\begin{tipbox}
最值得固化的，不是“所有参数”，而是你一周里会重复用很多次的那两三种工作模式。
\end{tipbox}

\begin{warnbox}
配置文件字段名和可用项可能会随版本演进变化。真正落地时，还是以你本地帮助和当前官方配置文档为准。先想清楚“我要固化哪几种模式”，再去写具体字段，通常更不容易乱。
\end{warnbox}

\subsection{非交互执行：什么时候该用 \texttt{codex exec}}

如果你想把某些工作做成脚本，或者在 CI 里跑，\texttt{codex exec} 很适合。比如：

\begin{lstlisting}[style=cli]
codex exec "总结当前仓库最重要的三个风险点"
\end{lstlisting}

它还支持把最终消息写到文件里：

\begin{lstlisting}[style=cli]
codex exec "生成本次改动总结" -o summary.txt
\end{lstlisting}

如果你要方便程序读取的输出，可以加：

\begin{lstlisting}[style=cli]
--json
\end{lstlisting}

如果你要严格结构化结果，可以再加：

\begin{lstlisting}[style=cli]
--output-schema
\end{lstlisting}

\subsection{CI 里怎么认证}

官方非交互文档明确建议：在 CI 里优先用 API 密钥，并通过 \texttt{CODEX\_API\_KEY} 注入。它还特别指出：\texttt{CODEX\_API\_KEY} 这个环境变量\textbf{只在 \texttt{codex exec} 里支持}。

\begin{lstlisting}[style=cli]
CODEX_API_KEY=<api-key> codex exec --json "梳理当前未处理的问题报告"
\end{lstlisting}

\subsection{把 \texttt{exec} 真正接进脚本：人看摘要，机器看结构}

很多人把 \texttt{codex exec} 用到“命令能跑通”就停了，但脚本闭环通常还差一步。更稳的自动化思路，是把输出拆成三层：
\begin{itemize}
  \item 给\textbf{人}看的最后摘要；
  \item 给\textbf{脚本}读的结构化事件流；
  \item 给\textbf{流程}判断的退出状态和文件存在性检查。
\end{itemize}

一个很实用的最小样板可以长这样：

\begin{lstlisting}[style=cli]
set -euo pipefail
export CODEX_API_KEY=your_api_key

codex exec --json \
  --output-schema review-schema.json \
  -o codex-summary.txt \
  "审查当前改动，并按约定结构输出风险摘要" \
  > codex-events.jsonl
\end{lstlisting}

这几个产物各有分工：
\begin{itemize}
  \item \texttt{codex-summary.txt}：给人看，一眼知道结论；
  \item \texttt{codex-events.jsonl}：给脚本留结构化事件流，后面要上传日志、做二次处理会方便很多；
  \item \texttt{review-schema.json}：约束最终输出格式，让下游流程更稳定。
\end{itemize}

跑完后，脚本里至少再补一个很小的检查：

\begin{lstlisting}[style=cli]
test -s codex-summary.txt
\end{lstlisting}

这样做的意义很直接：就算这轮任务没有正常产出摘要，流水线也能尽早失败，而不是把“空结果”继续往后传。

\begin{tipbox}
如果后面真的要自动消费结果，尽量让机器读 \texttt{--json} 或 \texttt{--output-schema} 约束过的输出，不要去硬拆自然语言段落。
\end{tipbox}

\subsection{为什么 \texttt{exec} 默认要求 Git 仓库}

官方文档也写得很清楚：\texttt{codex exec} 默认要求在 Git 仓库内运行，这是为了降低破坏性改动风险。如果你非常确定环境安全，才考虑：

\begin{lstlisting}[style=cli]
codex exec --skip-git-repo-check "..."
\end{lstlisting}

\begin{warnbox}
把一项任务做成自动化之前，先手工把它跑稳定。官方最佳实践也明确不建议过早把不稳定流程自动化。
\end{warnbox}

\begin{sumbox}
从“能聊天”到“能进工作流”，关键就三件事：看 diff、跑测试、把重复规则沉淀到 \texttt{AGENTS.md}。而脚本化场景则优先考虑 \texttt{codex exec}。
\end{sumbox}

\section{任务完成验收卡}

\begin{learnbox}
学完这一节，你会知道一轮任务做到哪一步才算真正收工，而不是“它说改完了，我就先信了”。
\end{learnbox}

很多人不是不会用 Codex，而是不知道什么时候该停。最稳的做法，是在结束前拿下面这张验收卡过一遍。

\rowcolors{2}{white}{SummaryLight}
\begin{longtable}{>{\raggedright\arraybackslash}p{0.22\textwidth} >{\raggedright\arraybackslash}p{0.30\textwidth} >{\raggedright\arraybackslash}p{0.34\textwidth}}
\toprule
检查项 & 最低标准 & 不够清楚时可以继续怎么问 \\
\midrule
\endhead
改动范围 & 只碰到预期文件，没顺手改一堆无关内容 & 请列出本轮改动文件，并说明每个文件为什么必须改 \\
验证情况 & 至少说明跑了哪些检查，或者明确哪些没跑 & 请只列出你实际执行过的验证，以及还没验证的部分 \\
结果说明 & 能用通俗的话讲清楚“这次到底改了什么” & 请用业务视角总结本次变化，不要只重复 diff \\
剩余风险 & 能指出最值得我再看的边界和潜在回归点 & 请列出当前改动最值得我复核的两个风险点 \\
提交准备 & 你自己能判断现在适不适合提交或继续补检查 & 基于当前状态，你建议现在提交，还是先补哪一步 \\
\bottomrule
\end{longtable}

如果你想让 Codex 按这张卡汇报，可以直接说：

\begin{lstlisting}[style=prompt]
请按“改动文件 / 验证情况 / 剩余风险 / 是否建议现在提交”四项总结这次工作。
\end{lstlisting}

\begin{warnbox}
不要把“它说自己验证过了”直接当验收结果。你真正要看的，是它到底验证了什么、没验证什么，以及这些检查够不够覆盖你的关键路径。
\end{warnbox}

\begin{sumbox}
任务完成的判断标准不是“它停下来了”，而是你已经看清改了什么、怎么验的、还剩什么风险，以及现在能不能安心提交。
\end{sumbox}

\section{常见错误与排查思路}

\begin{learnbox}
学完这一节，你遇到问题时至少知道先查什么，不会一报错就完全没方向。
\end{learnbox}

最常见的问题，通常不是模型不聪明，而是环境、目录、登录、权限、上下文这几类基础问题。

\subsection{排查思路图}

\begin{center}
\begin{tikzpicture}[
  q/.style={diamond, draw=Warn, fill=WarnLight, aspect=2, align=center, inner sep=2pt},
  r/.style={draw=black!55, rounded corners, fill=white, minimum width=3.7cm, minimum height=0.9cm, align=center},
  arr/.style={-Latex, thick, draw=Warn}
]
\node[q] (start) {出问题了};
\node[r, below=10mm of start] (dir) {先看当前目录对不对};
\node[r, below=7mm of dir] (cmd) {再看命令能不能识别};
\node[r, below left=9mm and 12mm of cmd] (auth) {检查登录与权限};
\node[r, below right=9mm and 12mm of cmd] (ctx) {检查上下文是否太宽或太少};
\node[r, below=15mm of $(auth)!0.5!(ctx)$] (retry) {缩小任务范围后重试};

\draw[arr] (start) -- (dir);
\draw[arr] (dir) -- (cmd);
\draw[arr] (cmd) -- (auth);
\draw[arr] (cmd) -- (ctx);
\draw[arr] (auth.south) -- ++(0,-5mm) -| ([xshift=-8mm]retry.north);
\draw[arr] (ctx.south) -- ++(0,-5mm) -| ([xshift=8mm]retry.north);
\end{tikzpicture}
\end{center}

\subsection{问题一：命令找不到}

\begin{lstlisting}[style=cli]
command not found: codex
\end{lstlisting}

优先检查：
\begin{itemize}
  \item 是否真的安装成功；
  \item 当前命令行环境能不能找到 \texttt{codex}；
  \item 终端会话是否需要重新打开。
\end{itemize}

\subsection{问题二：明明在项目里，却读不到对的内容}

先确认你是不是在正确目录：

\begin{lstlisting}[style=cli]
pwd
ls
\end{lstlisting}

如果目录没问题，再看你是不是把任务说得太宽了。

\subsection{问题三：回答太空泛}

大概率不是它完全不会，而是你的任务太宽。解决办法通常是：
\begin{itemize}
  \item 缩小目标；
  \item 指定文件、模块或功能；
  \item 明确输出格式，比如“列出文件路径 + 每个文件一句解释”。
\end{itemize}

\subsection{问题四：它总在问你要不要批准}

这不一定是坏事。多数时候只是当前审批策略较保守，或者它准备执行的命令本来就不属于安全读操作。

你可以先用：

\begin{lstlisting}[style=cli]
/status
/permissions
\end{lstlisting}

看当前会话权限，再决定要不要调模式。

\subsection{问题五：脚本里跑不通}

如果你在 CI 或脚本里用 \texttt{codex exec}，优先确认：
\begin{itemize}
  \item 认证是不是用的 \texttt{CODEX\_API\_KEY}；
  \item 当前目录是不是 Git 仓库；
  \item 是否需要 \texttt{--output-last-message} 或 \texttt{--json} 这类更适合自动化的输出方式。
\end{itemize}

\begin{warnbox}
高风险操作、删除操作、覆盖性修改，一定要自己复核。AI 可以帮你提速，但不能替你承担最终判断责任。
\end{warnbox}

\begin{sumbox}
排查时优先看五件事：目录、安装、登录、权限、任务范围。先把这五件事排清楚，通常已经能解决大部分问题。
\end{sumbox}

\section{使用建议与最佳实践}

\begin{learnbox}
学完这一节，你会知道怎样把 Codex CLI 用得更稳、更省时间，而不是越用越乱。
\end{learnbox}

\subsection{建议一：任务要具体}

尽量少说“帮我优化一下”，多说“帮我解释这个目录”“帮我定位登录逻辑”“只改文案，不改业务逻辑”。

\subsection{建议二：先小后大}

别一上来就扔超大任务。更稳妥的顺序通常是：
\begin{enumerate}
  \item 先解释；
  \item 再定位；
  \item 再计划；
  \item 最后再修改。
\end{enumerate}

\subsection{建议三：让它输出结构化结果}

比如直接要求：
\begin{itemize}
  \item 按文件列表输出；
  \item 每个文件一句话解释；
  \item 先结论，再细节；
  \item 最后给验证建议。
\end{itemize}

\subsection{建议四：让它知道怎么验证}

官方最佳实践明确提醒，不要让代理在不知道怎么构建、测试、做静态检查的前提下瞎试。所以最好在仓库里维护好 \texttt{AGENTS.md}，或者在提示中说明验证方式。

\subsection{建议五：复杂任务先计划}

当任务比较长、改动面比较广时，先用 \texttt{/plan} 或者直接让它先出计划，再决定是否实施。这样能明显减少“上来就改，结果越改越散”的问题。

\subsection{建议六：把重复规则沉淀下来}

如果你经常重复同样的话，就说明这条规则应该搬进 \texttt{AGENTS.md}、技能文件或配置，而不是每次重新输入。

\begin{tipbox}
把 Codex CLI 当成“会先打草稿、再陪你迭代”的工程搭档，比把它当成“全自动代打工具”稳定得多。
\end{tipbox}

\begin{sumbox}
最佳实践就一句话：任务说清楚，范围控住，验证补上，重复规则沉淀下来。
\end{sumbox}

\section{新手最容易踩的坑}

\begin{learnbox}
学完这一节，你会提前知道哪些地方最容易翻车，尤其适合第一次正式使用前快速扫一遍。
\end{learnbox}

\begin{enumerate}
  \item 没确认当前目录，就开始让它读代码或改文件。
  \item 任务描述太宽，结果很空，还误以为是工具不行。
  \item 一上来就给太大权限，自己却没想清楚风险。
  \item 看见它改完了就直接接受，没有自己看 diff 和跑验证。
  \item 把“解释、定位、修改、验证”四件事全揉进一个超长需求里。
  \item 在自动化场景里直接照搬交互式习惯，没有单独处理认证和输出。
  \item 有重复规则却不写进 \texttt{AGENTS.md}，导致每次都要重新解释。
\end{enumerate}

\begin{warnbox}
新手最大的问题通常不是不会敲命令，而是没有形成稳定节奏：先确认，再限制范围，再验证结果。
\end{warnbox}

\begin{sumbox}
避坑的关键不是记住一堆例外，而是养成流程：确认目录、收紧任务、看状态、审改动、做验证。
\end{sumbox}

\section{高风险任务边界清单}

\begin{learnbox}
学完这一节，你会更清楚哪些任务不适合一上来就让 Codex 放手执行，尤其是在你还没做人工复核的时候。
\end{learnbox}

高风险任务不是“完全不能让 Codex 参与”，而是\textbf{不能直接把执行权整包交出去}。下面这些情况，尤其值得先收紧范围、先看计划、先做人工确认。

\rowcolors{2}{white}{WarnLight}
\begin{longtable}{>{\raggedright\arraybackslash}p{0.30\textwidth} >{\raggedright\arraybackslash}p{0.40\textwidth}}
\toprule
这类任务为什么高风险 & 更稳的做法 \\
\midrule
\endhead
删除、覆盖、批量改名、批量格式化：一旦范围没控住，回滚和代码审查都会很痛苦 & 先要求它只列计划和影响范围，再按小批次执行，改完立刻看 \texttt{/diff} \\
破坏性 Git 操作：比如重写历史、强推、回退到旧提交 & 先自己决定策略，再让它辅助解释影响，不要把最终判断外包出去 \\
数据库迁移、数据回填、状态修复脚本：副作用大，出错后代价高 & 先让它写风险说明、回滚思路和验证步骤，再由你确认后执行 \\
密钥、认证缓存、生产配置：一旦泄漏或改错，后果往往比代码问题更严重 & 先把敏感文件和目录说清楚，必要时明确要求“只解释，不修改” \\
会触发外部副作用的命令：部署、发请求、写远端资源、修改系统状态 & 优先用只读方式分析，真正执行前一定要自己审一遍命令含义和影响范围 \\
\bottomrule
\end{longtable}

如果你已经知道这是高风险任务，一个很好用的开场方式是：

\begin{lstlisting}[style=prompt]
这是一个高风险任务。先不要执行，也不要改文件。
请先列出影响范围、主要风险、回滚方法，以及你建议我先看的最小变更范围。
\end{lstlisting}

\begin{tipbox}
高风险任务最怕的不是“工具不够聪明”，而是你没有先把边界、回滚和验证想清楚。先让它帮你分析，再决定要不要放行，通常更稳。
\end{tipbox}

\begin{sumbox}
判断一件事要不要谨慎放手，有个很简单的标准：如果做错之后难回滚、难做代码审查、代价高，那它就属于高风险任务。
\end{sumbox}

\section{进阶使用技巧：从会用到会控场}

\begin{learnbox}
学完这一节，你会开始把 Codex CLI 当成真正的工程搭档来用，不只是“帮我改一段代码”，而是会主动控范围、控节奏、控交付。
\end{learnbox}

前面那些内容已经够你顺利上手了。到了进阶阶段，重点就不再是“知道有哪些命令”，而是\textbf{知道什么时候该先想、什么时候该先收范围、什么时候该先比方案、什么时候该固定输出格式}。

\subsection{技巧一：先做方案评审，不急着改代码}

\textbf{适合什么时候：}需求还没完全定、同一个目标可能有两三种实现路径、你担心改动面会扩散的时候。\\
\textbf{核心思路：}先让 Codex 做“方案比较”，不是立刻进入“执行模式”。

\begin{lstlisting}[style=prompt]
先不要改代码。请给我 2 到 3 种实现方案，分别说明：
1. 需要改哪些文件
2. 改动范围大概多大
3. 主要优点和风险
4. 你更推荐哪一种，为什么
\end{lstlisting}

这类问法很适合放在需求刚成形的时候。它能帮你更早看到“最小改法”和“最顺手但改动更大”的区别。

\begin{warnbox}
常见误用是：嘴上说“先给方案”，后面又补一句“顺手直接改掉”。这样很容易让会话直接滑进执行阶段，失去先比较的价值。
\end{warnbox}

\subsection{技巧二：大仓库先控上下文，不要一口气全拉进来}

\textbf{适合什么时候：}仓库很大、目录很多、你已经开始感觉回答越来越空，或者会话里反复塞长日志、长代码的时候。\\
\textbf{核心思路：}先收文件范围，再收功能范围，最后再收动作范围。

\begin{lstlisting}[style=prompt]
先不要全仓库泛讲。请先找出和“登录失败提示”最相关的 5 个文件，
只输出文件路径，并说明我下一轮应该重点追哪两个。
\end{lstlisting}

这一招的关键，不是让它“知道更多”，而是让它\textbf{只盯真正相关的那一小块}。配合 \texttt{/mention} 和 \texttt{/compact}，通常会比不断贴更多背景更有效。

\begin{tipbox}
大仓库里一个很稳的顺序是：先问相关文件，再问调用链，再问最小改动点。先别一上来就要求“完整分析整个系统”。
\end{tipbox}

\subsection{技巧三：用 \texttt{/fork} 保留旧思路，同时试另一条线}

\textbf{适合什么时候：}你已经积累了一轮上下文，但想试另一种问题假设、另一种实现方案、或者另一种排查顺序。\\
\textbf{核心思路：}不要推翻当前会话，而是把它当成一条已经建立好的分支，再开一条线比较。

\begin{lstlisting}[style=cli]
/fork
\end{lstlisting}

你可以在分叉后的会话里直接说：

\begin{lstlisting}[style=prompt]
保留上一轮的上下文，但这次请不要沿用原来的结论。
请从“问题可能出在缓存或状态同步”这个角度重新给排查顺序。
\end{lstlisting}

\texttt{/fork} 最值钱的地方，是它能让你把“旧结论”和“新猜想”并排保留，而不是每次都从零重讲一遍。

\begin{warnbox}
常见误用是把 \texttt{/fork} 当“重开会话”。如果你只是想清掉当前聊天，不需要保留对比关系，那应该优先用 \texttt{/new} 或 \texttt{/clear}。
\end{warnbox}

\subsection{技巧四：用 \texttt{/model} 做快慢分工}

\textbf{适合什么时候：}你想先快速扫目录、列文件，再进入更复杂的分析、修改或代码审查。\\
\textbf{核心思路：}把“浏览定位”和“深入执行”拆开，不要整个任务都用同一种节奏。

\begin{lstlisting}[style=cli]
/model
\end{lstlisting}

一个很稳的工作节奏通常是：
\begin{enumerate}
  \item 先用更快的模型做扫目录、找入口、列文件；
  \item 再切到更强的模型做复杂推理、修改、代码审查；
  \item 如果后面又回到大量浏览工作，再切回更快的模型。
\end{enumerate}

这样做的价值不是“追求参数感”，而是让会话节奏更像真实工程工作：先摸清，再深挖。

\subsection{技巧五：先定义成功标准，再开始动手}

\textbf{适合什么时候：}任务边界模糊、你担心它顺手改太多、或者你自己也还没想清“做到什么算完成”。\\
\textbf{核心思路：}先把验收条件写出来，再让 Codex 开始做。

\begin{lstlisting}[style=prompt]
先不要修改。请先复述目标、列出不做的事、给出完成标准和验证方法。
等我确认后再开始。
\end{lstlisting}

这类写法特别像你在团队里先对齐需求。它能明显减少“它做了不少事，但不是我真正要的那件事”这种偏差。

\begin{tipbox}
如果你发现自己总在说“不是这个意思”“这个也别改”“这个先不做”，那就说明成功标准应该前置，而不是等它做完再修正。
\end{tipbox}

\subsection{技巧六：把代码审查做成闭环，不是只跑一次}

\textbf{适合什么时候：}你已经改完一轮，准备决定是否提交，或者你怀疑当前修改里还有行为风险。\\
\textbf{核心思路：}代码审查不是一句结束语，而是一轮小循环。

一个很实用的闭环是：
\begin{enumerate}
  \item 先修改；
  \item 跑一次 \texttt{/review}；
  \item 只优先修真正高风险问题；
  \item 再看 \texttt{/diff}；
  \item 最后做一轮简短复核。
\end{enumerate}

如果你想让它聚焦真正的问题，可以直接说：

\begin{lstlisting}[style=prompt]
请只列真正值得我在提交前处理的问题。
把“高风险问题”和“可选优化”分开写，先说高风险问题。
\end{lstlisting}

\begin{warnbox}
代码审查最常见的误用，是把“它给了一堆建议”当成“都得改”。真正应该先处理的，是那些会引起行为回归、边界错误、缺失验证的问题。
\end{warnbox}

\subsection{技巧七：固定结构化交付格式}

\textbf{适合什么时候：}你希望每轮结果都更像真实交付，而不是一大段自由发挥。\\
\textbf{核心思路：}不要每次临场听它发挥，直接固定输出模板。

\begin{lstlisting}[style=prompt]
请固定按下面格式汇报：
1. 结论
2. 改动文件
3. 验证情况
4. 剩余风险
5. 下一步建议
\end{lstlisting}

这招最适合和 \texttt{AGENTS.md} 一起用。因为一旦你发现团队每次都想看同一种结果格式，那它就应该从“临时提示词”升级成“仓库规则”。

\begin{tipbox}
结构化交付不是为了好看，而是为了让你更快判断：这轮到底改了什么、验了什么、还差什么、现在能不能交。
\end{tipbox}

\subsection{技巧八：把 \texttt{codex exec} 用在真正适合自动化的地方}

\textbf{适合什么时候：}你要做批量总结、风险归纳、分类整理、CI 检查、规范化输出，而不是大范围自动改代码。\\
\textbf{核心思路：}优先自动化“读和总结”，谨慎自动化“改和执行”。

\begin{lstlisting}[style=cli]
codex exec --json "总结当前分支相对 main 的三个主要风险"

codex exec --json \
  --output-schema schema.json \
  "把当前改动整理成结构化交付摘要"
\end{lstlisting}

进阶阶段，\texttt{exec} 最值得做的通常是这些事：
\begin{itemize}
  \item 生成结构化改动总结；
  \item 对当前 diff 做风险分类；
  \item 统一整理 issue、日志、报错；
  \item 在 CI 里做只读检查和摘要输出。
\end{itemize}

而下面这些事，通常不适合太早自动化：
\begin{itemize}
  \item 还没在本地跑稳的大流程；
  \item 改动范围大的自动修复；
  \item 带外部副作用的执行动作；
  \item 出错后很难回滚的操作。
\end{itemize}

\subsection{技巧九：先用一个小例子理解 MCP 到底值在哪}

\textbf{适合什么时候：}仓库外还有产品文档、内部知识库、工单系统、接口说明，或者别的外部工具数据，而你不想在提示词里手动复制来复制去。\\
\textbf{核心思路：}让 Codex 同时看到“本地代码”和“外部上下文”，但第一次一定先从\textbf{一个只读小场景}跑通。

先用命令确认你现在接了哪些 MCP 服务：

\begin{lstlisting}[style=cli]
codex mcp list
codex mcp add docs --url https://example.com/mcp
\end{lstlisting}

然后拿一个特别具体的任务试水。比如当前仓库里有退款逻辑，外部文档里写着退款接口规范，你就可以这样问：

\begin{lstlisting}[style=prompt]
请结合当前仓库里的退款实现，以及 docs MCP 里的“退款接口”文档，
分开告诉我：
1. 本地代码现在实际怎么做
2. 文档要求是什么
3. 两边有没有不一致
4. 如果我要先验证一处，最该从哪开始
\end{lstlisting}

这类场景里，MCP 的价值就会很清楚：
\begin{itemize}
  \item 只看本地代码，你知道“现在做成了什么样”；
  \item 只看文档，你知道“理论上应该怎么做”；
  \item 两边一起看，你才能更快找出偏差点和验证顺序。
\end{itemize}

第一次接 MCP，最值得先试的通常不是“让它直接调用很多外部系统”，而是下面这类只读任务：
\begin{itemize}
  \item 本地实现和外部文档是否一致；
  \item 仓库里的字段命名和接口定义是否对得上；
  \item 工单里描述的问题，能不能在代码里快速对上位置。
\end{itemize}

\begin{tipbox}
第一次接 MCP，不要一上来把工单、监控、文档、数据库全接上。先接一个只读来源，先做“对照、解释、定位”，跑顺后再往外扩。
\end{tipbox}

\begin{warnbox}
MCP 让上下文更丰富，不等于判断可以外包。特别是带写入能力的外部服务，依旧要沿用前面讲的审批、范围和回滚思路。
\end{warnbox}

\begin{sumbox}
进阶使用的核心，不是会更多命令，而是会控四件事：控范围、控节奏、控比较、控交付。你一旦开始主动控制这四件事，Codex 的稳定性会明显上一个台阶。
\end{sumbox}

\section{正文收束：3分钟速读总结}

到这里为止，正文主线已经讲完了。下面先用一页把前面的重点收口；再往后是速查和附录，适合按需跳读，不必再按顺序硬看。

\begin{tcolorbox}[colback=SummaryLight,colframe=Summary,title=如果你现在只想记住最重要的几点,fonttitle=\bfseries]
\begin{enumerate}
  \item 把 Codex CLI 当成终端里的代码搭档，而不是全自动代打工具。
  \item 安装后先做三件事：\texttt{codex --help}、登录、最小任务验证。
  \item 日常开发多用 \texttt{codex}，脚本化多用 \texttt{codex exec}；改完更稳的顺序是先 \texttt{/review}，再 \texttt{/diff}，最后自己验证。
  \item 本地开发优先从 \texttt{workspace-write} 这类中间档权限开始，不要急着全放开。
  \item 高质量提示词的关键是：目标、范围、限制、输出格式。
  \item 真正让 Codex 稳定好用的，不是一次神 prompt，而是 \texttt{AGENTS.md}、清晰测试方式和稳定工作流。
\end{enumerate}
\end{tcolorbox}

\section{常用命令速查表}

这一页不讲新知识，只把正文里出现过的高频入口收成一张表，方便你回头查。看完这页，下一页再配合“新手避坑清单”一起用，会更像一套随手能翻的工具卡。

\rowcolors{2}{white}{AccentLight}
\begin{longtable}{>{\ttfamily\raggedright\arraybackslash}p{0.26\textwidth} >{\raggedright\arraybackslash}p{0.20\textwidth} >{\raggedright\arraybackslash}p{0.38\textwidth}}
\toprule
命令/操作 & 用途 & 什么时候用 \\
\midrule
\endhead
codex & 进入交互式会话 & 日常本地开发、连续提问、逐步迭代时最常用 \\
codex --help & 查看帮助信息 & 忘了入口、参数、功能时先用 \\
codex --version & 查看版本 & 怀疑版本差异时先确认 \\
codex login & 发起登录 & 想显式登录而不是等首次运行触发 \\
codex login status & 查看登录状态 & 不确定当前是否已认证时 \\
codex exec "..." & 非交互执行 & 脚本化、CI、一次性任务 \\
codex review --uncommitted & 审查当前未提交改动 & 提交前快速做风险检查 \\
codex review --base main & 对比分支基线做评审 & 提交拉取请求（PR）前自查很实用 \\
codex resume --last & 恢复最近一次会话 & 不想重复讲背景时 \\
codex apply <TASK\_ID> & 应用某次任务产出的 diff & 你想把结果落到本地工作树时 \\
/status & 查看会话状态 & 确认模型、权限、可写目录 \\
/permissions & 调整权限模式 & 临时改成更保守或更自动 \\
/diff & 查看当前改动 & 改完先看一眼再提交 \\
/review & 在会话内做代码审查 & 想先查风险和缺失测试 \\
/plan & 先做执行计划 & 任务较大、较复杂时 \\
/init & 生成 \texttt{AGENTS.md} 脚手架 & 想把规则沉淀到仓库里 \\
\bottomrule
\end{longtable}

\section{新手避坑清单}

这一页和上一页速查表最好搭着看：上一页解决“现在该用哪个入口”，这一页解决“哪些地方最容易因为节奏不稳而翻车”。

\begin{itemize}
  \item 操作前先确认你在哪个目录。
  \item 登录方式和使用场景要区分：交互式和 CI 不是一回事。
  \item 能先解释就先解释，别急着大改。
  \item 让它改文件时，要求说明改了哪些文件、为什么改、怎么验证。
  \item 会话里看不对劲，先用 \texttt{/status} 和 \texttt{/permissions}。
  \item 改完更稳的顺序是先 \texttt{/review}，再看 \texttt{/diff}，最后自己验证。
  \item 有重复规则就写进 \texttt{AGENTS.md}，别每次重新说。
  \item 遇到版本差异时，以你本地 \texttt{codex --help} 和官方文档为准。
\end{itemize}

\section{附录：从零到一演练页}

\begin{learnbox}
学完这一节，你会看到一条从克隆公开小仓库到完成一次解释、定位、修改、复核的完整路径，适合第一次照着做。
\end{learnbox}

从这里开始进入附录区。前面的正文已经够你上手；后面的内容更偏演练、模板和查表。\\
这一页的重点不是仓库本身有多复杂，而是帮你把\textbf{真实工作流顺一遍}。它和前面的“最小可运行示例”也不重复：前者是最短体验，这一页是完整演练。这里用“任意公开小仓库”来演示，你可以换成一个最小待办应用、博客或练习项目。

\subsection{目标}

假设你的目标很简单：
\begin{itemize}
  \item 先看懂仓库结构；
  \item 再找到入口；
  \item 再做一个低风险改动；
  \item 最后复核并总结。
\end{itemize}

\subsection{第 1 步：把公开仓库拉到本地}

\begin{lstlisting}[style=cli]
git clone <公开小仓库地址> demo-repo
cd demo-repo
\end{lstlisting}

这里故意不把演示绑死到某一个具体仓库名上，是为了让你能把这页直接复用到任何公开小项目。

\subsection{第 2 步：第一次进入会话}

\begin{lstlisting}[style=cli]
codex
\end{lstlisting}

第一轮建议先这样问：

\begin{lstlisting}[style=prompt]
请先按新手视角查看当前项目。
告诉我：
1. 入口文件在哪里
2. 核心目录分别负责什么
3. 我应该按什么顺序阅读
请尽量引用具体文件路径。
\end{lstlisting}

\subsection{第 3 步：定位一个具体点}

等你拿到目录概览以后，继续把范围收紧：

\begin{lstlisting}[style=prompt]
现在请重点解释最核心的两个文件。我想知道它们分别负责什么，以及它们之间是怎么配合的。
\end{lstlisting}

\subsection{第 4 步：做一个低风险改动}

第一次练手，不要先碰大逻辑，优先选这类改动：
\begin{itemize}
  \item README 说明补充；
  \item 文案更清楚；
  \item 一个明显的小问题；
  \item 缺失但容易补的测试。
\end{itemize}

比如你可以这样说：

\begin{lstlisting}[style=prompt]
请只做一个低风险改动：把首页提示文案改得更清楚，不要改业务逻辑。修改前先告诉我你准备改哪些文件。
\end{lstlisting}

\subsection{第 5 步：复核改动}

改完以后不要立刻收工，先看：

\begin{lstlisting}[style=cli]
/review
/diff
\end{lstlisting}

再让它补一句人话总结：

\begin{lstlisting}[style=prompt]
请总结这次改动：
1. 改了哪些文件
2. 每个文件改了什么
3. 我现在应该怎么验证
\end{lstlisting}

\subsection{第 6 步：完成一轮最小闭环}

这时你应该已经完成了最重要的一轮：
\begin{enumerate}
  \item 进入仓库；
  \item 读懂结构；
  \item 定位关键点；
  \item 做一处低风险修改；
  \item 看 diff 和代码审查结果；
  \item 准备验证。
\end{enumerate}

\begin{tipbox}
第一次练手时，不追求改得多，追求完整走完一轮。能稳定走完这一轮，比一次改十个文件更重要。
\end{tipbox}

\begin{sumbox}
“从零到一演练页”的核心不是仓库名，而是让你掌握一条固定路径：先理解，再定位，再小改，再复核。
\end{sumbox}

\section{附录：练习题怎么看算合格}

\begin{learnbox}
学完这一节，你会更容易判断自己的练习结果算不算合格。这里不给“标准答案”，而是给“合格答案的样子”。
\end{learnbox}

因为每个仓库都不一样，所以这页不提供唯一答案。更有用的做法，是看你的结果有没有落到下面这些关键点。\\
这页故意放在“练习题页”前面，是因为你先知道判分标准，后面练的时候会更有尺子；如果你已经做完，也可以回来对照自检。

\subsection{练习一：合格结果至少要有路径和顺序}

合格结果通常至少要满足三件事：
\begin{itemize}
  \item 指出具体文件路径，而不是只说“入口大概在前端目录”；
  \item 说明为什么从这个文件开始看；
  \item 给出一个能真正照着走的阅读顺序。
\end{itemize}

\begin{lstlisting}[style=cli]
示例格式：
1. 入口文件：`src/main.ts`
2. 最值得先看的目录：`src/router/`、`src/pages/`
3. 推荐阅读顺序：
   - 先看 `src/main.ts`，确认应用怎么启动
   - 再看 `src/router/index.ts`，确认页面入口怎么串起来
   - 最后看 `src/pages/Login.vue`，把登录页逻辑补上
\end{lstlisting}

\subsection{练习二：合格结果要讲清调用链，不只是列文件}

如果你做的是“定位功能”练习，合格结果通常至少要满足三件事：
\begin{itemize}
  \item 能指出入口或触发点，而不是只堆文件名；
  \item 能说清最相关的 3 到 5 个文件怎么串起来；
  \item 能指出如果下一步只改一个小点，最可能先看哪里。
\end{itemize}

\begin{lstlisting}[style=cli]
示例格式：
1. 入口：`src/pages/Login.vue` 里的提交按钮触发登录动作
2. 关键文件：
   - `src/pages/Login.vue`：收集输入并触发表单提交
   - `src/store/auth.ts`：发起登录请求并更新登录状态
   - `src/api/auth.ts`：封装登录接口
3. 调用顺序：
   页面提交 -> store action -> API 请求 -> 成功后更新状态
4. 如果只改“登录失败提示文案”，先看 `src/pages/Login.vue`
\end{lstlisting}

\subsection{练习三：合格结果要控住范围，还要说清怎么验}

如果你做的是低风险小改动，合格结果通常会有这些特征：
\begin{itemize}
  \item 只改一两个真正相关的文件；
  \item 明确说出“没有改业务逻辑”或者“没有顺手重构”；
  \item 给出至少一条可执行的验证方法。
\end{itemize}

\begin{lstlisting}[style=cli]
示例格式：
- 改动文件：`src/pages/Login.vue`
- 改动原因：把登录失败提示改得更清楚
- 未改动部分：没有修改请求逻辑、校验逻辑、路由跳转
- 验证方法：输入错误密码，确认页面提示文案符合预期
\end{lstlisting}

\subsection{练习四和练习五：重点看它有没有真的“落地”}

\textbf{Review 练习} 合不合格，重点看这几件事：
\begin{itemize}
  \item 它有没有优先指出真实风险，而不是只夸；
  \item 它的评论有没有对上当前 diff；
  \item 它有没有把“必须改”和“可选优化”分开。
\end{itemize}

\textbf{提示词练习} 合不合格，重点看这几件事：
\begin{itemize}
  \item 改写后的提示词有没有明确目标、范围、限制、输出格式；
  \item 结果里是不是明显多了具体文件、动作和验证；
  \item 你是不是更容易继续追问，而不是继续被迫重开题。
\end{itemize}

\begin{tipbox}
自检时有个很实用的判断法：如果你的结果里能直接看到“路径、动作、限制、验证”，通常就已经比“空泛介绍型答案”强很多了。
\end{tipbox}

\begin{sumbox}
练习题最怕的不是答错，而是答得看起来很多，却没有路径、没有边界、没有验证。只要你能把这三样做出来，练习就已经很值了。
\end{sumbox}

\section{附录：好提示词 vs 坏提示词}

\begin{learnbox}
学完这一节，你会更容易判断什么样的提问会让 Codex 输出靠谱结果，什么样的提问容易把它带偏。
\end{learnbox}

\rowcolors{2}{white}{SummaryLight}
\begin{longtable}{p{0.24\textwidth} p{0.32\textwidth} p{0.32\textwidth}}
\toprule
场景 & 不太好的写法 & 更好的写法 \\
\midrule
\endhead
看项目 & 帮我看看这个项目 & 请按新手视角介绍当前项目，重点说入口、核心目录、主流程和推荐阅读顺序 \\
定位功能 & 帮我找登录 & 请帮我定位登录相关代码，告诉我入口、涉及文件和大致调用链 \\
处理报错 & 帮我修一下这个错 & 这是报错信息。请先解释最可能原因，再给排查顺序，最后给最小修改方案 \\
改代码 & 帮我优化一下 & 请只修改提示文案，不要改业务逻辑，修改后列出改动文件和验证方法 \\
做代码审查 & 帮我看看有没有问题 & 请像代码审查者一样检查当前改动，先列问题，再给简短总结，重点关注回归和缺失测试 \\
\bottomrule
\end{longtable}

\subsection{一个简单判断法}

如果一条提示词同时满足下面四件事，它通常就已经比平均水平强很多：
\begin{itemize}
  \item 说清楚目标；
  \item 说清楚范围；
  \item 说清楚限制；
  \item 说清楚输出格式。
\end{itemize}

\subsection{最常见的坏味道}

\begin{itemize}
  \item 只有动作，没有目标，比如“优化一下”“改一下”。
  \item 只有目标，没有范围，比如“把这个项目整理好”。
  \item 只有结果要求，没有约束，比如“修好它”，却没说能不能改架构。
  \item 没要求输出格式，导致结果全是大段空话。
\end{itemize}

\begin{sumbox}
真正的区别不在“高级词藻”，而在有没有把目标、范围、限制、输出格式交代清楚。
\end{sumbox}

\section{附录：常见场景小抄}

\begin{learnbox}
学完这一节，你可以把常见场景当成小抄来用，遇到任务时直接照着改，不用从头想。
\end{learnbox}

\subsection{场景一：接手项目}

\begin{tipbox}
先问什么：入口、核心目录、主流程、阅读顺序。\\
推荐提示词：请按新手视角介绍当前项目，引用具体文件路径。\\
别忘了：先让它解释，不要一上来就改。
\end{tipbox}

\subsection{场景二：处理报错}

\begin{tipbox}
先问什么：最可能原因、排查顺序、最值得看的文件。\\
推荐提示词：先解释报错，不要急着修改。\\
别忘了：让它给最小修改方案，而不是上来大改。
\end{tipbox}

\subsection{场景三：补测试}

\begin{tipbox}
先问什么：项目现有测试框架、最接近的已有测试、这次应该补哪种测试。\\
推荐提示词：请先找项目里现有测试模式，再补一条风格一致的测试。\\
别忘了：明确要求不要顺手重构无关代码。
\end{tipbox}

\subsection{场景四：做代码审查}

\begin{tipbox}
先问什么：行为风险、边界情况、缺失测试、可能回归。\\
推荐动作：先 \texttt{/review}，再 \texttt{/diff}。\\
别忘了：要求“先列问题，再总结”，这样输出更适合做代码审查。
\end{tipbox}

\begin{sumbox}
常见场景小抄的核心作用，是帮你减少临场组织语言的成本。越常见的任务，越值得模板化。
\end{sumbox}

\section{附录：CI 集成示例}

\begin{learnbox}
学完这一节，你会知道怎样把 \texttt{codex exec} 接进自动化流程，尤其是怎样组合 \texttt{--json}、\texttt{-o} 和 API 密钥。
\end{learnbox}

\subsection{最小命令行示例}

\begin{lstlisting}[style=cli]
export CODEX_API_KEY=your_api_key

codex exec --json \
  -o codex-summary.txt \
  "审查当前仓库，并总结最主要的 3 个风险点"
\end{lstlisting}

这一段里，三个关键点分别是：
\begin{itemize}
  \item \texttt{CODEX\_API\_KEY}：给非交互执行做认证。
  \item \texttt{--json}：让标准输出更适合被脚本读取。
  \item \texttt{-o codex-summary.txt}：把最后一条自然语言总结单独写到文件里。
\end{itemize}

\subsection{一个适合 CI 的思路}

你可以把 \texttt{codex exec} 放到下面这些位置：
\begin{itemize}
  \item 拉取请求（PR）创建后做风险总结；
  \item 每日巡检时总结仓库变化；
  \item 发布前做一轮只读式代码审查；
  \item 自动生成本次改动摘要。
\end{itemize}

\subsection{GitHub Actions 风格示例}

\begin{lstlisting}[style=cli]
name: codex-check

on:
  pull_request:

jobs:
  codex-review:
    runs-on: ubuntu-latest
    steps:
      - uses: actions/checkout@v4

      - name: Install Codex CLI
        run: npm i -g @openai/codex

      - name: Run Codex Exec
        env:
          CODEX_API_KEY: ${{ secrets.CODEX_API_KEY }}
        run: |
          codex exec --json -o codex-summary.txt \
            "审查当前改动相对基线分支的差异，并总结主要风险"
\end{lstlisting}

\subsection{CI 里最容易踩的坑}

\begin{itemize}
  \item 直接照搬交互式习惯，没有切到 \texttt{codex exec}。
  \item 忘了提供 API 密钥。
  \item 没有考虑当前目录是不是 Git 仓库。
  \item 只看 stdout，没有把最终摘要单独落文件。
\end{itemize}

\begin{warnbox}
自动化场景更适合做“总结、分类、只读式检查”这类任务。真正会改很多文件的流程，先在本地跑稳，再考虑自动化。
\end{warnbox}

\begin{sumbox}
CI 集成的关键不是把一切自动化，而是先找到最适合自动化的那一层：总结、审查、提示、归档。
\end{sumbox}

\section{附录：真实排错案例}

\begin{learnbox}
学完这一节，你会看到一条更接近真实开发现场的排错路径，不只是“报错了怎么办”，而是“遇到一个具体报错时，怎样让 Codex 帮你缩小范围”。
\end{learnbox}

这一页故意写成案例，而不是抽象原则。因为新手最容易卡住的点，不是“不知道要排错”，而是不知道\textbf{第一句该怎么问、第二步该看什么}。

\subsection{案例背景}

假设你在一个前端项目里运行开发命令时，终端出现类似错误：

\begin{lstlisting}[style=cli]
TypeError: Cannot read properties of undefined (reading 'map')
\end{lstlisting}

你目前只知道三件事：
\begin{itemize}
  \item 页面白屏了；
  \item 报错和某个列表渲染有关；
  \item 你还没定位到具体文件。
\end{itemize}

\subsection{第一轮，不要急着让它直接修}

先这样问：

\begin{lstlisting}[style=prompt]
我刚遇到一个运行时报错：
TypeError: Cannot read properties of undefined (reading 'map')

请先不要直接修改代码。
先帮我做三件事：
1. 解释这个错误最常见的原因
2. 告诉我应该先检查哪些类型的文件
3. 给我一个从快到慢的排查顺序
\end{lstlisting}

这一轮的目标不是立刻修复，而是让 Codex 先帮你把问题分类。通常它会给出这种方向：
\begin{itemize}
  \item 某个数组值还没准备好，就提前调用了 \texttt{map}；
  \item 接口返回为空，或者初始状态不是数组；
  \item 某个组件 props 没传到位。
\end{itemize}

\subsection{第二轮，把范围压到文件级别}

等第一轮拿到方向以后，再追问：

\begin{lstlisting}[style=prompt]
请根据当前项目代码，帮我定位最可能出现这个问题的文件。
优先找：
1. 渲染列表的组件
2. 初始化 state 的地方
3. 拉取接口数据并赋值的地方
\end{lstlisting}

这一步很关键。因为“抽象原因”到这里才真正开始落地到代码。

\subsection{第三轮，让它给最小修复方案}

等你已经知道大概是哪几个文件以后，再说：

\begin{lstlisting}[style=prompt]
现在请给我一个最小修复方案。
要求：
1. 先说明为什么会报错
2. 只改最必要的地方
3. 不要顺手重构
4. 修改后告诉我怎么验证
\end{lstlisting}

\subsection{这类案例里，Codex 最值得帮你的地方}

不是“替你秒修”，而是帮你更快完成下面这些动作：
\begin{itemize}
  \item 把报错归类；
  \item 找最可能的文件；
  \item 给出从小到大的排查顺序；
  \item 限制改动范围；
  \item 补验证建议。
\end{itemize}

\subsection{如果它一上来就想改很多文件怎么办}

立刻把范围收回来：

\begin{lstlisting}[style=prompt]
请先不要做大范围修改。只给我最小修复方案，并说明如果只改一个文件，最可能改哪一个。
\end{lstlisting}

\subsection{这页案例真正想教你的东西}

不是这个错误本身，而是一种顺序：
\begin{enumerate}
  \item 先解释报错；
  \item 再定位文件；
  \item 再最小修复；
  \item 最后验证。
\end{enumerate}

\begin{tipbox}
当你不知道该不该马上让 Codex 动代码时，默认先让它解释和定位。先把问题缩小，再让它改，几乎总是更稳。
\end{tipbox}

\begin{sumbox}
真实排错案例的核心，不是记住某一种报错，而是记住排错顺序：先分类，后定位，再最小修复，最后验证。
\end{sumbox}

\section{附录：练习题页}

\begin{learnbox}
学完这一节，你可以把前面学到的用法真正练一遍，而不是只停留在“我看懂了”。
\end{learnbox}

这一页给你几道适合新手的练习题。最好找一个公开小仓库，或者你自己的练手仓库，按顺序做。\\
如果你还没看上一页“练习题怎么看算合格”，建议先扫一眼；做题时你会更清楚什么叫真的落地。

\subsection{练习一：找到阅读入口}

\textbf{目标}：让 Codex 帮你给出项目阅读顺序。\\
\textbf{要求}：
\begin{itemize}
  \item 问出入口文件；
  \item 问出核心目录；
  \item 问出推荐阅读顺序；
  \item 输出里必须包含具体文件路径。
\end{itemize}

\textbf{自检点}：
\begin{itemize}
  \item 结果有没有落到具体路径；
  \item 结果是不是按新手视角组织的；
  \item 你能不能根据它给的顺序自己开始读代码。
\end{itemize}

\subsection{练习二：定位一个功能}

\textbf{目标}：找出某个功能的入口和调用链。\\
\textbf{要求}：
\begin{itemize}
  \item 自己先选一个功能，比如登录、上传、搜索、列表渲染；
  \item 让 Codex 告诉你涉及哪些文件；
  \item 让它解释调用链，不要只列文件名。
\end{itemize}

\textbf{自检点}：
\begin{itemize}
  \item 它有没有解释“文件之间怎么串起来”；
  \item 你能不能据此自己继续往下追。
\end{itemize}

\subsection{练习三：做一个低风险小改动}

\textbf{目标}：让 Codex 改一个你敢接受的小地方。\\
\textbf{要求}：
\begin{itemize}
  \item 只改文案、说明或一处很小的逻辑；
  \item 明确要求“不要改业务逻辑”或“不要顺手重构”；
  \item 修改后让它列出改动文件和验证方法。
\end{itemize}

\textbf{自检点}：
\begin{itemize}
  \item 改动范围有没有被控制住；
  \item 它有没有解释为什么这么改；
  \item 你有没有自己看 \texttt{/diff}。
\end{itemize}

\subsection{练习四：做一次代码审查}

\textbf{目标}：把代码审查流程走一遍。\\
\textbf{建议动作}：

\begin{lstlisting}[style=cli]
/review
/diff
\end{lstlisting}

\textbf{自检点}：
\begin{itemize}
  \item 它有没有优先指出真正的问题，而不是只给表扬；
  \item 你能不能分清“高风险问题”和“可选优化”；
  \item 你有没有把它的审查结果和实际 diff 对照起来看。
\end{itemize}

\subsection{练习五：写一条更好的提示词}

\textbf{目标}：练习把模糊需求改成清晰需求。\\
\textbf{做法}：
\begin{itemize}
  \item 先写一句模糊请求，比如“帮我优化一下这里”；
  \item 再自己把它改写成包含目标、范围、限制、输出格式的版本；
  \item 然后分别交给 Codex，看结果差异。
\end{itemize}

\textbf{自检点}：
\begin{itemize}
  \item 改写后的结果是不是明显更落地；
  \item 你是不是更容易继续追问；
  \item 结果里是不是少了空话，多了具体文件和动作。
\end{itemize}

\begin{tipbox}
练习题不要一次做太多。对新手来说，最有效的顺序通常是：练习一、练习二、练习三。先把“理解、定位、小改动”这三件事练熟。
\end{tipbox}

\begin{sumbox}
练习题页的重点不是难度，而是帮助你把教程里的方法变成自己的操作习惯。
\end{sumbox}

\section{附录：AGENTS.md 示例}

\begin{learnbox}
学完这一节，你会看到一份可以直接改成自己仓库版本的 \texttt{AGENTS.md} 样板，也会知道为什么这个文件对长期使用 Codex 很重要。
\end{learnbox}

前面已经提到，\texttt{AGENTS.md} 的作用是把重复规则沉淀下来。你可以把它理解成“这个仓库专属的长期系统提示词”。这样一来，Codex 每次进入仓库时都更容易知道：
\begin{itemize}
  \item 项目怎么跑；
  \item 测试怎么跑；
  \item 哪些目录能改，哪些别碰；
  \item 交付结果时你最看重什么。
\end{itemize}

\subsection{什么样的内容最值得写进去}

最值得沉淀的，通常是下面这些信息：
\begin{itemize}
  \item 仓库结构和重点目录说明；
  \item 安装、启动、测试、静态检查的标准命令；
  \item 改动边界，比如哪些文件是生成的、哪些目录不要动；
  \item 代码风格、提交前检查、输出格式要求；
  \item 完成任务后应该如何总结和汇报。
\end{itemize}

\subsection{一份适合中小型项目起步的样板}

\begin{lstlisting}[style=cli]
# AGENTS.md

## 项目概览

- 这个仓库包含前端、后端和共享工具模块。
- 改动的首要目标是：范围小、容易评审、方便验证。
- 在做较大改动之前，优先先阅读现有代码和测试。

## 仓库结构

- `frontend/`：界面和页面代码
- `backend/`：接口处理和业务逻辑
- `shared/`：前后端共用代码
- `scripts/`：本地维护脚本
- `tests/`：自动化测试
- `docs/`：项目文档

## 工作规则

- 修改前先确认当前目录和相关文件。
- 优先做小而明确的改动，不要轻易做大范围重构。
- 除非任务明确要求，否则不要修改生成文件。
- 除非确有必要，否则不要重命名文件或移动模块。
- 如果需要高风险或破坏性改动，先说明原因再执行。

## 常用命令

- 安装依赖：`npm install`
- 启动开发环境：`npm run dev`
- 运行单元测试：`npm test`
- 运行静态检查：`npm run lint`
- 构建项目：`npm run build`

## 改动原则

- 保持现有代码风格和项目约定。
- 不要顺手修改无关文件。
- 只有在代码本身不够直观时才添加注释。
- 处理问题时优先选择最小且正确的改动。

## 验证要求

- 改完代码后，先执行最小但相关的验证步骤。
- 如果任务影响了应用行为，尽量运行相关测试。
- 如果有必须做但无法运行的检查，请在最终总结里明确说明。

## 输出要求

- 用通俗的话总结这次改了什么。
- 如果任务不算特别简单，请列出改动过的文件。
- 说明已经做了哪些验证。
- 如果还有剩余风险或后续检查项，也请写出来。

## 默认策略

- 遇到模糊需求时，先澄清问题或先给一个简短计划。
- 遇到较大任务时，先拆成更小的步骤。
- 如果用户还在熟悉代码库，优先先解释，再修改。
\end{lstlisting}

\subsection{如果你的仓库更复杂，可以再加什么}

如果是团队项目，还可以继续往里补：
\begin{itemize}
  \item 分支策略，比如不要直接动主分支；
  \item 数据库迁移规范；
  \item 前后端联调顺序；
  \item 哪些命令耗时长，不要默认跑；
  \item 哪些目录包含敏感信息或外部生成产物。
\end{itemize}

\subsection{写 AGENTS.md 的一个实用原则}

不要把它写成“项目说明书大全”。最实用的 \texttt{AGENTS.md}，通常只保留\textbf{Codex 真的会反复用到的规则}。如果一条说明不会影响它怎么读代码、怎么改代码、怎么验证结果，那就不一定要放进来。

\begin{tipbox}
一个好用的判断标准是：如果你在三次任务里都重复说过同一句话，那它大概率就该进 \texttt{AGENTS.md} 了。
\end{tipbox}

\begin{sumbox}
\texttt{AGENTS.md} 的核心作用，不是展示项目背景，而是把“怎么做事”固定下来。写得越贴近真实工作流，Codex 后面越省心。
\end{sumbox}

\section{附录：版本变更提示框}

\begin{learnbox}
学完这一节，你会知道当教程、截图和你本地命令不一致时，应该优先检查哪些地方。
\end{learnbox}

\begin{versionbox}
\begin{itemize}
  \item 安装方式可能变化：例如 npm 包名、推荐安装命令、首次登录入口。
  \item 登录方式可能变化：浏览器登录、设备码登录、API 密钥登录的推荐路径可能调整。
  \item 斜杠命令可能变化：命令名、输出内容和支持范围都有可能随版本演进。
  \item \texttt{exec} 和 \texttt{review} 的参数可能增加或调整。
  \item 沙箱和审批策略的默认行为可能变化，尤其是新版本加入新策略时。
  \item 最可靠的检查方式，始终是你本地的 \texttt{codex --help} 和对应子命令帮助。
\end{itemize}
\end{versionbox}

\subsection{最推荐的检查顺序}

\begin{enumerate}
  \item 先看 \texttt{codex --help}；
  \item 再看 \texttt{codex <subcommand> --help}；
  \item 再回到官方文档页面查对应章节；
  \item 最后再参考博客、截图、旧教程。
\end{enumerate}

\subsection{哪些地方最容易发生版本差异}

\begin{itemize}
  \item 首次登录流程；
  \item 斜杠命令列表；
  \item \texttt{codex exec} 的自动化参数；
  \item \texttt{review} 的输出样式；
  \item 审批与沙箱策略的默认组合。
\end{itemize}

\begin{sumbox}
当版本变化发生时，不要先怀疑自己记错了，也不要先怀疑教程错了。先看本地 help，这是最快的判断方法。
\end{sumbox}

\section{FAQ}

\begin{learnbox}
学完这一节，你会对几个最常见的新手问题有直接答案，尤其是“上下文窗口超了怎么办”这种很容易卡住的问题。
\end{learnbox}

\subsection{Q1：如果上下文窗口超了，应该怎么办？}

最实用的做法不是“硬塞更多内容”，而是\textbf{主动压缩上下文}。可以按这个顺序处理：
\begin{enumerate}
  \item 先缩小任务范围，只围绕一个功能、一个报错或一组文件继续问。
  \item 用 \texttt{/compact} 压缩长会话，把已经确认的背景收成短摘要。
  \item 用 \texttt{/mention} 或直接点名文件路径，只把真正关键的文件拉进当前上下文。
  \item 如果当前话题已经太散，直接开一轮新会话，并把上一轮结论用 3 到 5 句话重新总结后贴进去。
\end{enumerate}

\begin{tipbox}
判断标准很简单：如果你已经开始反复贴长日志、长代码、长历史对话，就说明该收缩上下文了。先总结，再继续，通常比硬接着聊更稳。
\end{tipbox}

\subsection{Q2：它回答得很空泛，是模型不行吗？}

大多数时候不是。更常见的原因是任务太宽、目标不够具体，或者没有要求它按文件、按步骤、按输出格式回答。先把问题缩到“哪个功能、哪几个文件、最后要什么格式”，效果通常会马上变好。

\subsection{Q3：什么时候该开新会话，而不是一直接着问？}

当你发现当前会话已经同时混进了多个目标，比如“看项目结构、处理报错、改文案、做 CI”，这时就该拆。一个会话尽量只服务一个主目标，后续会更稳，也更容易复盘。

\subsection{Q4：会话到底怎么切？能像别的工具那样在多个会话间切来切去吗？}

可以，但要分清楚是哪一种“切”。

\begin{itemize}
  \item 想在当前终端里立刻开一轮新的：用 \texttt{/new}
  \item 想清屏后再开新的：用 \texttt{/clear}
  \item 想切回某个旧会话：用 \texttt{/resume}
  \item 想保留旧会话内容，但分一条新线：用 \texttt{/fork}
\end{itemize}

如果你期待的是像某些工具那样，在同一界面里长期挂着多个会话标签页来回切换，那 Codex CLI 的主思路不是这个。它更像是：\textbf{把之前的会话记录起来，需要时再恢复、分叉或新开。}

\subsection{Q5：文档里是不是还有命令没讲？}

现在这份长版里已经补上了\textbf{完整命令参考}，包括顶层命令、常见子命令，以及高频全局参数怎么读。\\
如果你只是第一次上手，还是建议先走主线章节；如果你准备查命令全景、补边角知识、对照本地版本，直接去后面的“附录：完整命令参考”会更高效。

\subsection{Q6：它改完代码后，我最少要做什么检查？}

如果你只做最低限度检查，最少做三件事：
\begin{enumerate}
  \item 看 \texttt{/diff}，确认改动范围对不对；
  \item 跑最相关的测试、构建或静态检查；
  \item 自己走一遍最关键的用户路径。
\end{enumerate}

如果时间允许，再补一次 \texttt{/review} 会更稳，因为它更容易提前指出行为风险和漏掉的测试。

\begin{sumbox}
FAQ 里最值得先记的一条就是：上下文一旦开始发散，就收缩任务、压缩会话、必要时新开一轮。不要指望靠堆更多上下文来解决所有问题。
\end{sumbox}

\section{附录：完整命令参考}

\begin{learnbox}
学完这一节，你会对当前版本 Codex CLI 的顶层命令有完整地图，也知道每个命令大概该在什么场景下用。你不用把它全背下来，但至少查的时候不会迷路。
\end{learnbox}

这一节根据你本机 \texttt{codex --help} 以及对应子命令帮助整理。它的目标不是把帮助输出原样抄一遍，而是把“\textbf{命令是干什么的、什么时候值得用、哪些参数最关键}”讲清楚。\\
另一个容易混淆的点也先说清楚：\textbf{这一节讲的是 CLI 顶层命令和子命令。}会话里的斜杠命令，比如 \texttt{/status}、\texttt{/review}、\texttt{/compact}，前面已经单独讲过了。

\subsection{先看顶层命令总览}

\rowcolors{2}{white}{SummaryLight}
\begin{longtable}{>{\ttfamily\raggedright\arraybackslash}p{0.19\textwidth} >{\raggedright\arraybackslash}p{0.25\textwidth} >{\raggedright\arraybackslash}p{0.42\textwidth}}
\toprule
命令 & 它是干什么的 & 你大概什么时候会用到 \\
\midrule
\endhead
codex [PROMPT] & 启动交互式会话 & 日常本地开发、连续追问、边看边改时最常用 \\
codex exec & 非交互执行代理任务 & 脚本、CI、一次性任务、结构化输出 \\
codex review & 非交互代码审查 & 想直接对未提交改动、分支差异或某个 commit 做代码审查 \\
codex login & 管理登录 & 首次登录、切换认证方式、检查登录状态 \\
codex logout & 清掉本地认证 & 需要登出、换账号、清理本机凭据时 \\
codex mcp & 管理外部 MCP 服务 & 给 Codex 接外部工具、外部数据源或第三方服务时 \\
codex mcp-server & 把 Codex 自己作为 MCP 服务暴露出去 & 做更底层的集成、桥接或工具链对接时 \\
codex app-server & 启动或辅助应用服务端相关工具 & 更偏集成、协议生成、IDE 或远端应用服务端场景 \\
codex completion & 生成命令行补全脚本 & 想让 Bash / Zsh / Fish 里补全更顺手时 \\
codex sandbox & 手动在 Codex 提供的沙箱里跑命令 & 想测试沙箱行为、单独验证受限执行环境时 \\
codex debug & 调试工具集合 & 排查应用服务端或更底层行为时 \\
codex apply & 把任务产出的 diff 应用到本地工作树 & 你已经有任务结果，只想把补丁应用到仓库里时 \\
codex resume & 恢复旧交互会话 & 不想重新讲背景，想接着上次继续干 \\
codex fork & 从旧会话分叉出一条新线 & 想保留原结论，再试另一种思路时 \\
codex cloud & 浏览或提交 Codex Cloud 任务 & 用云端任务池、云端执行、再把结果落回本地时 \\
codex features & 查看或切换功能开关 & 想检查某个功能开关是否启用，或手动开关实验特性 \\
codex help & 查看帮助 & 忘了命令、参数、子命令时最快的入口 \\
\bottomrule
\end{longtable}

\begin{tipbox}
如果你现在只想把主线抓住，先盯住 8 个：\texttt{codex}、\texttt{exec}、\texttt{review}、\texttt{login}、\texttt{logout}、\texttt{apply}、\texttt{resume}、\texttt{fork}。剩下那些命令，更多是扩展、集成或调试层面的能力。
\end{tipbox}

\subsection{第一组：主线工作流命令}

\textbf{\texttt{codex [PROMPT]}}：这是默认入口。你不写子命令时，Codex 会直接进入交互式会话。\\
它最适合“边问边追、边看边改”的工作方式。你甚至可以把第一句需求直接作为参数带进去：

\begin{lstlisting}[style=cli]
codex
codex "先帮我梳理当前仓库入口、核心目录和推荐阅读顺序"
\end{lstlisting}

这条命令支持最多的“会话级参数”，比如选模型、选沙箱、选审批策略、切配置档、加搜索、加额外可写目录等。如果帮助里说这些选项会转发给交互式 CLI，意思就是这些参数大多会作用在这一轮交互式 CLI 里。

\medskip
\textbf{\texttt{codex exec [PROMPT] [COMMAND]}}：这是非交互入口。它适合脚本、CI、批量总结、结构化输出，不适合那种需要你来回追问很多轮的任务。\\
它当前还有两个子命令：
\begin{itemize}
  \item \texttt{codex exec resume}：在非交互模式下恢复旧会话，并追加一段提示词。
  \item \texttt{codex exec review}：在非交互模式下做代码审查，适合脚本化调用。
\end{itemize}

最常见的几类参数是：
\begin{itemize}
  \item \texttt{--json}：把事件流输出成 JSONL，方便程序读取；
  \item \texttt{-o} / \texttt{--output-last-message}：把最后一条消息单独写到文件；
  \item \texttt{--output-schema}：要求最终输出符合某个 JSON Schema；
  \item \texttt{--skip-git-repo-check}：允许不在 Git 仓库里运行；
  \item \texttt{--ephemeral}：本轮不把会话文件持久化到磁盘。
\end{itemize}

你可以这样用：

\begin{lstlisting}[style=cli]
codex exec --json -o summary.txt "总结当前改动最主要的三个风险点"
\end{lstlisting}

\medskip
\textbf{\texttt{codex review [PROMPT]}}：这是非交互代码审查入口。你可以把它理解成“直接跑一轮审查模式”，不用先进入交互界面。\\
最关键的三个定位参数是：
\begin{itemize}
  \item \texttt{--uncommitted}：看当前已暂存、未暂存和未跟踪的改动；
  \item \texttt{--base <BRANCH>}：拿当前改动和某个基线分支对比；
  \item \texttt{--commit <SHA>}：只审查某个 commit 引入的改动。
\end{itemize}

\begin{lstlisting}[style=cli]
codex review --uncommitted
codex review --base main
\end{lstlisting}

如果你要在脚本里读取结果，又想保留代码审查这层语义，也可以考虑 \texttt{codex exec review}，因为它额外支持 \texttt{--json}、\texttt{-o}、\texttt{--ephemeral} 这类非交互参数。

\medskip
\textbf{\texttt{codex login}}：管理登录。它本身最重要的两个场景，一个是第一次登录，一个是切换认证方式。\\
当前最值得知道的入口有三个：
\begin{itemize}
  \item \texttt{codex login}：走默认登录流程；
  \item \texttt{codex login status}：检查当前登录状态；
  \item \texttt{codex login --with-api-key}：从标准输入读 API 密钥；
  \item \texttt{codex login --device-auth}：适合远端 Linux、SSH 或无浏览器环境。
\end{itemize}

\begin{lstlisting}[style=cli]
codex login
codex login status
printenv OPENAI_API_KEY | codex login --with-api-key
\end{lstlisting}

\medskip
\textbf{\texttt{codex logout}}：清掉当前本地认证信息。它的用途很直白，就是登出。\\
最常见的使用场景是：
\begin{itemize}
  \item 你准备换账号；
  \item 你想清理本机认证缓存；
  \item 你怀疑本地认证状态出了问题，想重新登录。
\end{itemize}

\medskip
\textbf{\texttt{codex apply <TASK\_ID>}}：把某个任务产出的 diff 应用到你本地工作树。帮助里说得很直接，它底层会把结果作为 \texttt{git apply} 打到工作树里。\\
这条命令特别适合“结果已经有了，但我现在只想把补丁应用到本地”这种场景。你需要给它一个 \texttt{TASK\_ID}，所以它不是日常第一次上手就会高频用到的入口。

\medskip
\textbf{\texttt{codex resume [SESSION\_ID] [PROMPT]}}：恢复旧交互会话。默认会打开一个选择列表；如果你只想接着最近一次，直接加 \texttt{--last}。\\
它还支持：
\begin{itemize}
  \item \texttt{--all}：显示全部会话，而不是只看当前工作目录附近的；
  \item \texttt{--include-non-interactive}：把非交互会话也纳入候选；
  \item 交互会话那批常见参数，比如 \texttt{-m}、\texttt{-s}、\texttt{-a}、\texttt{--search}、\texttt{--add-dir}。
\end{itemize}

\begin{lstlisting}[style=cli]
codex resume --last
codex resume 123e4567-e89b-12d3-a456-426614174000 "继续上次的排查，重点看缓存层"
\end{lstlisting}

\medskip
\textbf{\texttt{codex fork [SESSION\_ID] [PROMPT]}}：从旧会话分叉。它和 \texttt{resume} 的差别不是“能不能回到旧上下文”，而是“保不保留原线继续存在”。\\
当你想把旧结论留着，同时试另一个假设、另一个方案、另一个排障方向时，用它最顺。

\begin{lstlisting}[style=cli]
codex fork --last
codex fork 123e4567-e89b-12d3-a456-426614174000 "这次改从状态同步问题的角度再看一遍"
\end{lstlisting}

\medskip
\textbf{\texttt{codex help [COMMAND]}}：最朴素，但也最可靠。你忘了什么，就查什么。\\
如果你只记一条“不会过时”的建议，那就是：

\begin{lstlisting}[style=cli]
codex help
codex help exec
codex exec --help
\end{lstlisting}

\subsection{第二组：扩展、集成和云端命令}

\textbf{\texttt{codex mcp}}：管理外部 MCP（Model Context Protocol）服务。你可以把它理解成“给 Codex 接外部能力”的入口。\\
当前子命令有 6 个，分工很清楚：

\rowcolors{2}{white}{AccentLight}
\begin{longtable}{>{\ttfamily\raggedright\arraybackslash}p{0.20\textwidth} >{\raggedright\arraybackslash}p{0.52\textwidth}}
\toprule
子命令 & 用途 \\
\midrule
\endhead
mcp list & 查看当前配置过的 MCP 服务；加 \texttt{--json} 可以输出方便程序读取的结果 \\
mcp get <NAME> & 看某个 MCP 服务的具体配置；也支持 \texttt{--json} \\
mcp add <NAME> --url <URL> & 添加一个基于 HTTP 的 MCP 服务 \\
mcp add <NAME> -- <COMMAND>... & 添加一个通过 stdio 启动的 MCP 服务，可配 \texttt{--env KEY=VALUE} \\
mcp remove <NAME> & 删除一个 MCP 服务配置 \\
mcp login <NAME> & 对某个 MCP 服务做 OAuth 登录；可加 \texttt{--scopes} \\
mcp logout <NAME> & 清掉某个 MCP 服务的认证状态 \\
\bottomrule
\end{longtable}

\begin{lstlisting}[style=cli]
codex mcp list
codex mcp add my-server --url https://example.com/mcp
codex mcp add local-tools -- /path/to/mcp-server --flag
\end{lstlisting}

\medskip
\textbf{\texttt{codex mcp-server}}：把 Codex 自己作为 MCP 服务启动起来，当前帮助里写的是 stdio 模式。\\
这条命令对普通日常开发者不算高频，更适合做工具链集成、桥接、IDE 插件、或让别的客户端把 Codex 当服务来调。

\medskip
\textbf{\texttt{codex app-server}}：实验性命令。它既能启动应用服务端，也能做协议相关生成。\\
顶层最关键的是：
\begin{itemize}
  \item \texttt{--listen}：选择监听端点，默认是 \texttt{stdio://}，也可以走 websocket；
  \item 一组 \texttt{--ws-*} 参数：控制 websocket 鉴权方式、token 文件、签名配置等；
  \item \texttt{generate-ts}：生成 TypeScript 绑定；
  \item \texttt{generate-json-schema}：生成 JSON Schema。
\end{itemize}

\begin{lstlisting}[style=cli]
codex app-server --listen stdio://
codex app-server generate-ts --out ./generated
codex app-server generate-json-schema --out ./schema
\end{lstlisting}

\medskip
\textbf{\texttt{codex cloud}}：实验性云端任务入口。它的定位不是本地交互，而是“看云端任务、提云端任务、把结果拉回来”。\\
当前子命令有 5 个：

\rowcolors{2}{white}{SummaryLight}
\begin{longtable}{>{\ttfamily\raggedright\arraybackslash}p{0.22\textwidth} >{\raggedright\arraybackslash}p{0.50\textwidth}}
\toprule
子命令 & 用途 \\
\midrule
\endhead
cloud exec --env <ENV\_ID> [QUERY] & 提交一个新的 Cloud 任务；可加 \texttt{--attempts} 和 \texttt{--branch} \\
cloud status <TASK\_ID> & 看某个 Cloud 任务的状态 \\
cloud list & 列任务；支持 \texttt{--env}、\texttt{--limit}、\texttt{--cursor}、\texttt{--json} \\
cloud diff <TASK\_ID> & 看某个任务的 unified diff；可指定 \texttt{--attempt} \\
cloud apply <TASK\_ID> & 把某个任务的 diff 应用到本地；也可指定 \texttt{--attempt} \\
\bottomrule
\end{longtable}

这条命令最像“云端版的任务队列”。如果你当前主要是在本地终端里工作，可以先知道有它，不一定马上会高频使用。

\medskip
\textbf{\texttt{codex features}}：查看或切换功能开关。它更像一个“功能开关面板”，不是日常主线命令。\\
当前子命令是：
\begin{itemize}
  \item \texttt{features list}：看已知功能开关以及当前启用状态；
  \item \texttt{features enable <FEATURE>}：在 \texttt{config.toml} 里开启一个功能；
  \item \texttt{features disable <FEATURE>}：在 \texttt{config.toml} 里关闭一个功能。
\end{itemize}

\begin{lstlisting}[style=cli]
codex features list
codex features enable unified_exec
codex features disable unified_exec
\end{lstlisting}

\subsection{第三组：环境、沙箱和调试命令}

\textbf{\texttt{codex completion [SHELL]}}：生成命令行补全脚本。帮助里支持的命令行环境包括 \texttt{bash}、\texttt{zsh}、\texttt{fish}、\texttt{powershell}、\texttt{elvish}。\\
这条命令本身不直接改你的命令行配置，它只是把补全脚本打印出来或生成出来，后续通常要你自己接到对应配置里。

\begin{lstlisting}[style=cli]
codex completion bash
codex completion zsh
\end{lstlisting}

\medskip
\textbf{\texttt{codex sandbox}}：手动在 Codex 提供的沙箱里跑命令。这个命令不是“让代理在沙箱里工作”的入口，而是你自己显式去调用沙箱能力。\\
当前有 3 个平台子命令：

\rowcolors{2}{white}{AccentLight}
\begin{longtable}{>{\ttfamily\raggedright\arraybackslash}p{0.24\textwidth} >{\raggedright\arraybackslash}p{0.48\textwidth}}
\toprule
子命令 & 用途 \\
\midrule
\endhead
sandbox linux [COMMAND]... & 在 Linux 沙箱里执行命令；支持 \texttt{--full-auto} \\
sandbox macos [COMMAND]... & 在 macOS Seatbelt 下执行；额外支持 \texttt{--log-denials} \\
sandbox windows [COMMAND]... & 在 Windows restricted token 沙箱下执行 \\
\bottomrule
\end{longtable}

因为你现在主要写 Linux 教程，所以最值得记的是：

\begin{lstlisting}[style=cli]
codex sandbox linux --full-auto bash -lc 'pwd && ls'
\end{lstlisting}

它更适合“验证沙箱行为本身”，而不是替代你平时直接运行 \texttt{codex}。

\medskip
\textbf{\texttt{codex debug}}：调试工具集合。它当前主要围绕应用服务端调试。\\
当前层级是：
\begin{itemize}
  \item \texttt{debug app-server}：应用服务端调试工具入口；
  \item \texttt{debug app-server send-message-v2 <USER\_MESSAGE>}：向调试目标发一条消息。
\end{itemize}

这类命令更偏内部调试、协议调试、集成排障。对绝大多数“我只是想写代码、改文档、做代码审查”的使用者来说，不属于第一层必学内容。

\subsection{高频全局参数怎么读}

很多帮助里你会反复看到同一批参数。它们不是“每个命令都支持”，但确实是一批\textbf{高频复用参数}。先认熟这些，帮助阅读速度会快很多。

\subsubsection*{会话类命令常见参数}

\rowcolors{2}{white}{SummaryLight}
\begin{longtable}{>{\ttfamily\raggedright\arraybackslash}p{0.28\textwidth} >{\raggedright\arraybackslash}p{0.44\textwidth}}
\toprule
参数 & 作用 \\
\midrule
\endhead
-c, --config <key=value> & 临时覆盖 \texttt{\textasciitilde/.codex/config.toml} 的某个配置项 \\
--enable / --disable <FEATURE> & 临时打开或关闭某个功能开关 \\
-m, --model <MODEL> & 指定本轮会话要用的模型 \\
--oss & 改走本地开源模型提供方，而不是默认的云端提供方 \\
--local-provider <OSS\_PROVIDER> & 细分本地模型提供方，比如 \texttt{lmstudio} 或 \texttt{ollama} \\
-p, --profile <CONFIG\_PROFILE> & 从配置文件里选一个配置档 \\
-s, --sandbox <SANDBOX\_MODE> & 选沙箱级别：\texttt{read-only}、\texttt{workspace-write}、\texttt{danger-full-access} \\
-a, --ask-for-approval <APPROVAL\_POLICY> & 选审批策略，比如 \texttt{untrusted}、\texttt{on-request}、\texttt{never} \\
--full-auto & 一个比较省事的默认组合，通常等于更自动的审批 + \texttt{workspace-write} \\
\shortstack[l]{\texttt{--dangerously-bypass-}\\\texttt{approvals-and-sandbox}} & 完全跳过审批和沙箱，风险极高，只适合你自己非常清楚边界时 \\
-C, --cd <DIR> & 指定工作根目录 \\
-i, --image <FILE> & 把图片作为初始输入附上 \\
--search & 给模型打开联网搜索能力 \\
--add-dir <DIR> & 在主工作目录之外，再额外加一个可写目录 \\
--no-alt-screen & 关闭备用屏模式，保留终端滚动历史 \\
\bottomrule
\end{longtable}

\subsubsection*{非交互和自动化高频参数}

\rowcolors{2}{white}{AccentLight}
\begin{longtable}{>{\ttfamily\raggedright\arraybackslash}p{0.28\textwidth} >{\raggedright\arraybackslash}p{0.44\textwidth}}
\toprule
参数 & 作用 \\
\midrule
\endhead
--skip-git-repo-check & 允许不在 Git 仓库里运行 \\
--ephemeral & 不把会话文件持久化到磁盘 \\
--json & 把事件或结果输出成 JSONL，方便脚本处理 \\
-o, --output-last-message <FILE> & 把最后一条消息单独写到文件 \\
--output-schema <FILE> & 要求最终输出符合某个 JSON Schema \\
--uncommitted & 审查当前未提交改动 \\
--base <BRANCH> & 审查当前改动相对某个基线分支的变化 \\
--commit <SHA> & 审查某个 commit 引入的变化 \\
--title <TITLE> & 给审查摘要一个标题 \\
\bottomrule
\end{longtable}

\begin{warnbox}
帮助里经常会重复出现 \texttt{-c}、\texttt{--enable}、\texttt{--disable}。它们不是“这条命令有什么特别功能”，而是整套 CLI 共用的一种配置入口。真正决定场景差异的，通常还是命令本身和那条命令独有的参数。
\end{warnbox}

\subsection{最后给一个查命令的顺序}

当你忘了某个命令具体怎么用，最稳的顺序通常是：
\begin{enumerate}
  \item 先看这节，判断它属于哪一类命令；
  \item 再跑 \texttt{codex help <command>} 或 \texttt{codex <command> --help}；
  \item 如果它还有子命令，再继续看 \texttt{codex <command> <subcommand> --help}；
  \item 最后再结合你当前项目场景决定要不要真的执行。
\end{enumerate}

\begin{sumbox}
真正的“完整命令参考”不是把帮助背下来，而是知道三件事：这条命令在哪一层、适合什么场景、我下一步该查哪一级 help。只要这三件事清楚了，你基本就不会在命令树里迷路。
\end{sumbox}

\section{官方资源与版本说明}

\begin{learnbox}
学完这一节，你会知道这份教程背后的主要参考来源，也知道当本地行为和教程不一致时应该信谁。
\end{learnbox}

这份教程的关键事实主要来自两类来源：
\begin{itemize}
  \item 本地 CLI 帮助输出：\texttt{codex --help}、\texttt{codex exec --help}、\texttt{codex review --help}、\texttt{codex login --help} 等。
  \item OpenAI 官方 Codex 文档。
\end{itemize}

建议优先查看这些官方页面：
\begin{itemize}
  \item CLI 总览：\url{https://developers.openai.com/codex/cli}
  \item 认证：\url{https://developers.openai.com/codex/auth}
  \item 审批与安全：\url{https://developers.openai.com/codex/agent-approvals-security}
  \item 命令行选项：\url{https://developers.openai.com/codex/cli/reference}
  \item 斜杠命令：\url{https://developers.openai.com/codex/cli/slash-commands}
  \item 非交互模式：\url{https://developers.openai.com/codex/noninteractive}
  \item 最佳实践：\url{https://developers.openai.com/codex/learn/best-practices}
\end{itemize}

\subsection{当教程和你本地行为不一致时怎么办}

优先级建议这样排：
\begin{enumerate}
  \item 先看你本地的 \texttt{codex --help}；
  \item 再看对应的官方文档页面；
  \item 最后再看第三方博客、帖子或旧截图。
\end{enumerate}

\begin{advancedbox}
原因很简单：Codex CLI 更新比较快，而第三方教程经常会落后于当前版本。
\end{advancedbox}

\begin{sumbox}
真正可依赖的来源，始终是“你机器上的帮助信息 + OpenAI 官方 Codex 文档”。这两者要优先于任何二手资料。
\end{sumbox}

\newpage
\section{一页式总结页}

\begin{center}
\begin{tikzpicture}[
  card/.style={draw=Accent, rounded corners, fill=AccentLight, text width=5.5cm, minimum height=2.35cm, align=left, inner sep=6pt}
]
\node[card] (c1) at (-3.6,0) {\textbf{1. 先做什么}\\进入项目目录，确认位置没错。};
\node[card] (c2) at (3.6,0) {\textbf{2. 先跑什么}\\先看 \texttt{codex --help}，再登录，再做最小验证。};
\node[card] (c3) at (-3.6,-3.25) {\textbf{3. 日常怎么用}\\交互式用 \texttt{codex}，边问边收紧范围。};
\node[card] (c4) at (3.6,-3.25) {\textbf{4. 改完怎么查}\\先 \texttt{/review}，再 \texttt{/diff}，再自己验证。};
\node[card] (c5) at (-3.6,-6.5) {\textbf{5. 权限怎么选}\\本地优先 \texttt{workspace-write}，别急着全放开。};
\node[card] (c6) at (3.6,-6.5) {\textbf{6. 一句话记忆}\\说清楚任务，小步推进，看状态，审改动，做验证。};
\end{tikzpicture}
\end{center}

\end{document}
